\documentclass[t]{beamer}

% Load general definitions
% Preamble file - general definitions, package loading, etc.

%=================================
% Load packages
\usepackage{amssymb,amsmath}
\usepackage{graphicx}
\usepackage{url}
\usepackage{tikz}
\usetikzlibrary{mindmap,trees,arrows}
\usepackage{fancyvrb}
\usepackage[portuguese]{babel} 
\usepackage[utf8]{inputenc}
\usepackage{subfigure}
\usepackage{times}
\usepackage[T1]{fontenc}
\usepackage{cancel}
\usepackage{color}
\usepackage{listings}
\usepackage[document]{ragged2e}

%=================================
% Set mode
\mode<presentation>
{
	\usetheme{Madrid}
	\usecolortheme{structure}
	\useoutertheme{infolines}
	\setbeamercovered{invisible}
}

% Get rid of nav bar
\beamertemplatenavigationsymbolsempty

% Insert frame number at bottom of the page.
\usefoottemplate{\hfil\tiny{\color{black!90}\insertframenumber}} 

%=================================
% Define new commands

\newcommand\Real{{\mathbb{R}}}
%\newcommand{\vi}{\vspace{0.6\baselineskip}}
%\newcommand{\goodgap}{\hspace{\subfigtopskip}\hspace{\subfigbottomskip}}


% Equation environments
\newcommand{\beq}{\begin{equation}}
\newcommand{\eq}{\end{equation}}
\newcommand{\beqs}{\begin{equation*}}
\newcommand{\eqs}{\end{equation*}}
\newcommand{\beqn}{\begin{eqnarray}}
\newcommand{\eqn}{\end{eqnarray}}
% Bold variables
\newcommand{\mbf}[1]{\ensuremath{\mathbf{#1}}}
% Itemization
\newcommand{\bitem}{\begin{itemize}}
\newcommand{\eitem}{\end{itemize}}
\newcommand{\spitem}{\vskip 1em\item}
\newcommand{\bitems}{\begin{itemize}\item}
\newcommand{\benums}{\begin{enumerate}\item}
\newcommand{\eenum}{\end{enumerate}}
% color blocks
\newenvironment{colorblock}[2]{%
\setbeamercolor{block title}{#2}
\begin{block}{#1}}{\end{block}}
% Vertical spacing
\newcommand{\vone}{\vskip 1em}
\newcommand{\vhalf}{\vskip .5em}
% Frame environments
\newenvironment{ftst}[3][t]{%
\begin{frame}{environment=ftst,#1}
\frametitle{#2}
\framesubtitle{#3}}{\end{frame}}
\newenvironment{ftstf}[2]{
\begin{frame}[fragile,environment=ftstf]
\frametitle{#1}
\framesubtitle{#2}}{\end{frame}}
% colors
\definecolor{MyGray}{rgb}{0.5,0.5,0.5}
\definecolor{MyDBGray}{rgb}{0.1,0.1,0.4}
\definecolor{darkgreen}{rgb}{0,0.4,0}
\definecolor{black}{rgb}{0,0,0}
\def\defn#1{{\color{red} #1}}
% Footnote
\renewcommand{\thefootnote}{\alph{footnote}}
% Relaxed footnotes
\newcommand{\lfr}[1]{\let\thefootnote\relax\footnote{\tiny #1}}
% Verbatim environment - using FANCYVRB package
\DefineVerbatimEnvironment%
{rcode}{Verbatim}
{fontsize=\scriptsize}
% Verbatim environment - using LISTINGS package
%\lstnewenvironment{rcode} {\lstset{	language = R,
%									basicstyle = \scriptsize\ttfamily,
%									showspaces = false,
%									showstringspaces = false,
%									showtabs = false,
%									keywordstyle = \color{black}\bfseries,
%									commentstyle = \color{darkgreen},
%									numbers = none,
%									otherkeywords={	<-,
%													ggplot,
%													geom_boxplot,
%													facet_grid,
%													shapiro.test,
%													fligner.test,
%													glht,
%													with},
%									deletekeywords={data,
%													model,
%													residuals,
%													c,
%													axis,
%													default,
%													labels,
%													qq.text}}}%
%{}

% Specific definitions
\title[]{Algoritmos e Estruturas de Dados III}
\subtitle[]{Tipos de Grafos - Parte 2}
\author[]{Patrícia Lucas\\{\footnotesize }}
\institute{Bacharelado em Sistemas de Informação \\ IFNMG  - Campus Salinas}
\date{\scriptsize Salinas\\Dezembro 2020}

\begin{document}

% cover page
\setbeamertemplate{footline}{}
\begin{frame}

\begin{center}
\includegraphics[width=.15\textwidth]{}
\end{center}
  \titlepage
  \begin{tikzpicture}[remember picture,overlay]
  \node[anchor=south east,xshift=-5pt,yshift=5pt] at (current page.south east) {\tiny Versão 1.2020};
  \node[anchor=south west,yshift=0pt] at (current page.south west) {\includegraphics[width=.25\textwidth]{Logos/salinas_horizontal_jpg.jpg}};
  \end{tikzpicture}  
\end{frame}

% Main slides
\begin{ftst}{Grafo bipartido}{Tipos de Grafos}

\begin{itemize}
    \item Um grafo é bipartido se seus vértices podem ser divididos em dois conjuntos disjuntos não-vazios $N_1$ e $N_2$, tal que dois vértices são adjacentes se, e somente se, um deles pertence a $N_1$ e outro a $N_2$.
    \item Se cada vértice de $N_1$ é adjacente a todos os vértices de $N_2$, e vice-versa, temos um grafo bipartido completo. 
    \item Notação do exemplo: $K_{5,3}$
\end{itemize}
\vone
\centering


\tikzset{every picture/.style={line width=0.75pt}} %set default line width to 0.75pt        

\begin{tikzpicture}[x=0.75pt,y=0.75pt,yscale=-1,xscale=1]
%uncomment if require: \path (0,154); %set diagram left start at 0, and has height of 154

%Straight Lines [id:da08787424755786755] 
\draw [color={rgb, 255:red, 121; green, 120; blue, 120 }  ,draw opacity=1 ][line width=1.5]    (26.33,25.67) -- (93.33,111.67) ;
%Straight Lines [id:da8366360637059884] 
\draw [color={rgb, 255:red, 121; green, 120; blue, 120 }  ,draw opacity=1 ][line width=1.5]    (26.33,25.67) -- (161.33,111.67) ;
%Straight Lines [id:da9176534569392747] 
\draw [color={rgb, 255:red, 121; green, 120; blue, 120 }  ,draw opacity=1 ][line width=1.5]    (26.33,25.67) -- (230.33,111.67) ;
%Straight Lines [id:da9673799593150498] 
\draw [color={rgb, 255:red, 121; green, 120; blue, 120 }  ,draw opacity=1 ][line width=1.5]    (94.33,25.67) -- (93.33,111.67) ;
%Straight Lines [id:da9177957931776033] 
\draw [color={rgb, 255:red, 121; green, 120; blue, 120 }  ,draw opacity=1 ][line width=1.5]    (94.33,25.67) -- (161.33,111.67) ;
%Straight Lines [id:da5389732162626228] 
\draw [color={rgb, 255:red, 121; green, 120; blue, 120 }  ,draw opacity=1 ][line width=1.5]    (94.33,25.67) -- (230.33,111.67) ;
%Straight Lines [id:da46655353671150035] 
\draw [color={rgb, 255:red, 121; green, 120; blue, 120 }  ,draw opacity=1 ][line width=1.5]    (161.33,26.67) -- (93.33,111.67) ;
%Straight Lines [id:da4489760611616307] 
\draw [color={rgb, 255:red, 121; green, 120; blue, 120 }  ,draw opacity=1 ][line width=1.5]    (161.33,26.67) -- (161.33,111.67) ;
%Straight Lines [id:da20884936645803154] 
\draw [color={rgb, 255:red, 121; green, 120; blue, 120 }  ,draw opacity=1 ][line width=1.5]    (161.33,26.67) -- (230.33,111.67) ;
%Straight Lines [id:da6805945546525516] 
\draw [color={rgb, 255:red, 121; green, 120; blue, 120 }  ,draw opacity=1 ][line width=1.5]    (229.33,26.67) -- (93.33,111.67) ;
%Straight Lines [id:da5947235819521057] 
\draw [color={rgb, 255:red, 121; green, 120; blue, 120 }  ,draw opacity=1 ][line width=1.5]    (229.33,26.67) -- (161.33,111.67) ;
%Straight Lines [id:da9899362238635772] 
\draw [color={rgb, 255:red, 121; green, 120; blue, 120 }  ,draw opacity=1 ][line width=1.5]    (229.33,26.67) -- (230.33,111.67) ;
%Straight Lines [id:da9676387426333821] 
\draw [color={rgb, 255:red, 121; green, 120; blue, 120 }  ,draw opacity=1 ][line width=1.5]    (298.33,26.67) -- (93.33,111.67) ;
%Straight Lines [id:da5551081244310967] 
\draw [color={rgb, 255:red, 121; green, 120; blue, 120 }  ,draw opacity=1 ][line width=1.5]    (298.33,26.67) -- (161.33,111.67) ;
%Straight Lines [id:da8111281530603562] 
\draw [color={rgb, 255:red, 121; green, 120; blue, 120 }  ,draw opacity=1 ][line width=1.5]    (298.33,26.67) -- (230.33,111.67) ;
%Shape: Ellipse [id:dp419382005090444] 
\draw  [draw opacity=0][fill={rgb, 255:red, 245; green, 166; blue, 35 }  ,fill opacity=1 ] (18.92,25.67) .. controls (18.92,21.66) and (22.24,18.41) .. (26.33,18.41) .. controls (30.42,18.41) and (33.74,21.66) .. (33.74,25.67) .. controls (33.74,29.69) and (30.42,32.94) .. (26.33,32.94) .. controls (22.24,32.94) and (18.92,29.69) .. (18.92,25.67) -- cycle ;
%Shape: Ellipse [id:dp13527889449765773] 
\draw  [draw opacity=0][fill={rgb, 255:red, 245; green, 166; blue, 35 }  ,fill opacity=1 ] (86.92,25.67) .. controls (86.92,21.66) and (90.24,18.41) .. (94.33,18.41) .. controls (98.42,18.41) and (101.74,21.66) .. (101.74,25.67) .. controls (101.74,29.69) and (98.42,32.94) .. (94.33,32.94) .. controls (90.24,32.94) and (86.92,29.69) .. (86.92,25.67) -- cycle ;
%Shape: Ellipse [id:dp9827434474550794] 
\draw  [draw opacity=0][fill={rgb, 255:red, 245; green, 166; blue, 35 }  ,fill opacity=1 ] (153.92,26.67) .. controls (153.92,22.66) and (157.24,19.41) .. (161.33,19.41) .. controls (165.42,19.41) and (168.74,22.66) .. (168.74,26.67) .. controls (168.74,30.69) and (165.42,33.94) .. (161.33,33.94) .. controls (157.24,33.94) and (153.92,30.69) .. (153.92,26.67) -- cycle ;
%Shape: Ellipse [id:dp6487909355011356] 
\draw  [draw opacity=0][fill={rgb, 255:red, 245; green, 166; blue, 35 }  ,fill opacity=1 ] (221.92,26.67) .. controls (221.92,22.66) and (225.24,19.41) .. (229.33,19.41) .. controls (233.42,19.41) and (236.74,22.66) .. (236.74,26.67) .. controls (236.74,30.69) and (233.42,33.94) .. (229.33,33.94) .. controls (225.24,33.94) and (221.92,30.69) .. (221.92,26.67) -- cycle ;
%Shape: Ellipse [id:dp24712277624446166] 
\draw  [draw opacity=0][fill={rgb, 255:red, 245; green, 166; blue, 35 }  ,fill opacity=1 ] (290.92,26.67) .. controls (290.92,22.66) and (294.24,19.41) .. (298.33,19.41) .. controls (302.42,19.41) and (305.74,22.66) .. (305.74,26.67) .. controls (305.74,30.69) and (302.42,33.94) .. (298.33,33.94) .. controls (294.24,33.94) and (290.92,30.69) .. (290.92,26.67) -- cycle ;
%Shape: Ellipse [id:dp2212658359147779] 
\draw  [draw opacity=0][fill={rgb, 255:red, 15; green, 118; blue, 237 }  ,fill opacity=1 ] (85.92,111.67) .. controls (85.92,107.66) and (89.24,104.41) .. (93.33,104.41) .. controls (97.42,104.41) and (100.74,107.66) .. (100.74,111.67) .. controls (100.74,115.69) and (97.42,118.94) .. (93.33,118.94) .. controls (89.24,118.94) and (85.92,115.69) .. (85.92,111.67) -- cycle ;
%Shape: Ellipse [id:dp45346143959200624] 
\draw  [draw opacity=0][fill={rgb, 255:red, 15; green, 118; blue, 237 }  ,fill opacity=1 ] (153.92,111.67) .. controls (153.92,107.66) and (157.24,104.41) .. (161.33,104.41) .. controls (165.42,104.41) and (168.74,107.66) .. (168.74,111.67) .. controls (168.74,115.69) and (165.42,118.94) .. (161.33,118.94) .. controls (157.24,118.94) and (153.92,115.69) .. (153.92,111.67) -- cycle ;
%Shape: Ellipse [id:dp33011628323636644] 
\draw  [draw opacity=0][fill={rgb, 255:red, 15; green, 118; blue, 237 }  ,fill opacity=1 ] (222.92,111.67) .. controls (222.92,107.66) and (226.24,104.41) .. (230.33,104.41) .. controls (234.42,104.41) and (237.74,107.66) .. (237.74,111.67) .. controls (237.74,115.69) and (234.42,118.94) .. (230.33,118.94) .. controls (226.24,118.94) and (222.92,115.69) .. (222.92,111.67) -- cycle ;
%Shape: Rectangle [id:dp9906985720194883] 
\draw  [color={rgb, 255:red, 247; green, 42; blue, 42 }  ,draw opacity=1 ][dash pattern={on 1.69pt off 2.76pt}][line width=1.5]  (5,10) -- (321.11,10) -- (321.11,44.26) -- (5,44.26) -- cycle ;
%Shape: Rectangle [id:dp23577206552973107] 
\draw  [color={rgb, 255:red, 247; green, 42; blue, 42 }  ,draw opacity=1 ][dash pattern={on 1.69pt off 2.76pt}][line width=1.5]  (4.28,93.26) -- (320.38,93.26) -- (320.38,126.41) -- (4.28,126.41) -- cycle ;

% Text Node
\draw (331.5,21.39) node [anchor=north west][inner sep=0.75pt]    {$N_{1}$};
% Text Node
\draw (331.5,102.39) node [anchor=north west][inner sep=0.75pt]    {$N_{2}$};


\end{tikzpicture}

\end{ftst}

%=====

\begin{ftst}{Grafo bipartido}{Tipos de Grafos}

Um grafo G é bipartido se e somente se todo ciclo de G possuir comprimento par.
\vone

\centering 

\tikzset{every picture/.style={line width=0.75pt}} %set default line width to 0.75pt        

\begin{tikzpicture}[x=0.75pt,y=0.75pt,yscale=-1,xscale=1]
%uncomment if require: \path (0,166); %set diagram left start at 0, and has height of 166

%Shape: Ellipse [id:dp30212074499598573] 
\draw  [draw opacity=0][fill={rgb, 255:red, 15; green, 118; blue, 237 }  ,fill opacity=1 ] (22.91,21.87) .. controls (25.76,19.05) and (30.41,19.13) .. (33.28,22.04) .. controls (36.16,24.95) and (36.17,29.6) .. (33.32,32.42) .. controls (30.46,35.24) and (25.82,35.16) .. (22.94,32.25) .. controls (20.07,29.34) and (20.05,24.69) .. (22.91,21.87) -- cycle ;
%Shape: Ellipse [id:dp20331843151300388] 
\draw  [draw opacity=0][fill={rgb, 255:red, 15; green, 118; blue, 237 }  ,fill opacity=1 ] (22.52,92.56) .. controls (25.37,89.74) and (30.02,89.82) .. (32.89,92.73) .. controls (35.77,95.64) and (35.78,100.29) .. (32.93,103.11) .. controls (30.07,105.93) and (25.43,105.85) .. (22.55,102.94) .. controls (19.68,100.03) and (19.66,95.38) .. (22.52,92.56) -- cycle ;
%Shape: Ellipse [id:dp5989074558213376] 
\draw  [draw opacity=0][fill={rgb, 255:red, 15; green, 118; blue, 237 }  ,fill opacity=1 ] (93.6,22.43) .. controls (96.46,19.61) and (101.1,19.68) .. (103.98,22.59) .. controls (106.85,25.5) and (106.87,30.15) .. (104.01,32.97) .. controls (101.16,35.79) and (96.51,35.72) .. (93.63,32.8) .. controls (90.76,29.89) and (90.74,25.25) .. (93.6,22.43) -- cycle ;
%Shape: Ellipse [id:dp5951974638807662] 
\draw  [draw opacity=0][fill={rgb, 255:red, 15; green, 118; blue, 237 }  ,fill opacity=1 ] (93.09,92.82) .. controls (95.95,90) and (100.59,90.07) .. (103.47,92.99) .. controls (106.34,95.9) and (106.36,100.54) .. (103.5,103.36) .. controls (100.65,106.18) and (96,106.11) .. (93.13,103.2) .. controls (90.25,100.29) and (90.24,95.64) .. (93.09,92.82) -- cycle ;
%Straight Lines [id:da43106489597522435] 
\draw [color={rgb, 255:red, 16; green, 121; blue, 243 }  ,draw opacity=1 ][fill={rgb, 255:red, 0; green, 0; blue, 0 }  ,fill opacity=1 ][line width=1.5]    (27.72,97.84) -- (98.3,98.09) ;
%Straight Lines [id:da12872312542296527] 
\draw [color={rgb, 255:red, 16; green, 121; blue, 243 }  ,draw opacity=1 ][fill={rgb, 255:red, 0; green, 0; blue, 0 }  ,fill opacity=1 ][line width=1.5]    (98.3,98.09) -- (98.81,27.7) ;
%Straight Lines [id:da7485744883507535] 
\draw [color={rgb, 255:red, 16; green, 121; blue, 243 }  ,draw opacity=1 ][fill={rgb, 255:red, 0; green, 0; blue, 0 }  ,fill opacity=1 ][line width=1.5]    (28.11,27.15) -- (27.72,97.84) ;
%Straight Lines [id:da7016306132121852] 
\draw [color={rgb, 255:red, 16; green, 121; blue, 243 }  ,draw opacity=1 ][fill={rgb, 255:red, 0; green, 0; blue, 0 }  ,fill opacity=1 ][line width=1.5]    (28.11,27.15) -- (98.81,27.7) ;
%Shape: Ellipse [id:dp4475881908742718] 
\draw  [draw opacity=0][fill={rgb, 255:red, 15; green, 118; blue, 237 }  ,fill opacity=1 ] (57.52,57.56) .. controls (60.37,54.74) and (65.02,54.82) .. (67.89,57.73) .. controls (70.77,60.64) and (70.78,65.29) .. (67.93,68.11) .. controls (65.07,70.93) and (60.43,70.85) .. (57.55,67.94) .. controls (54.68,65.03) and (54.66,60.38) .. (57.52,57.56) -- cycle ;
%Shape: Ellipse [id:dp5795820585883926] 
\draw  [draw opacity=0][fill={rgb, 255:red, 15; green, 118; blue, 237 }  ,fill opacity=1 ] (129.09,57.82) .. controls (131.95,55) and (136.59,55.07) .. (139.47,57.99) .. controls (142.34,60.9) and (142.36,65.54) .. (139.5,68.36) .. controls (136.65,71.18) and (132,71.11) .. (129.13,68.2) .. controls (126.25,65.29) and (126.24,60.64) .. (129.09,57.82) -- cycle ;
%Straight Lines [id:da10605119277341335] 
\draw [color={rgb, 255:red, 16; green, 121; blue, 243 }  ,draw opacity=1 ][fill={rgb, 255:red, 0; green, 0; blue, 0 }  ,fill opacity=1 ][line width=1.5]    (63.2,62.62) -- (133.9,63.17) ;
%Straight Lines [id:da9560198451889184] 
\draw [color={rgb, 255:red, 16; green, 121; blue, 243 }  ,draw opacity=1 ][fill={rgb, 255:red, 0; green, 0; blue, 0 }  ,fill opacity=1 ][line width=1.5]    (28.11,27.15) -- (98.3,98.09) ;
%Straight Lines [id:da8268067794602478] 
\draw [color={rgb, 255:red, 16; green, 121; blue, 243 }  ,draw opacity=1 ][fill={rgb, 255:red, 0; green, 0; blue, 0 }  ,fill opacity=1 ][line width=1.5]    (27.72,97.84) -- (134.3,63.09) ;
%Straight Lines [id:da3311367596854673] 
\draw [color={rgb, 255:red, 16; green, 121; blue, 243 }  ,draw opacity=1 ][fill={rgb, 255:red, 0; green, 0; blue, 0 }  ,fill opacity=1 ][line width=1.5]    (98.81,27.7) -- (134.3,63.09) ;
%Shape: Ellipse [id:dp04087013508423709] 
\draw  [draw opacity=0][fill={rgb, 255:red, 15; green, 118; blue, 237 }  ,fill opacity=1 ] (202.6,22.43) .. controls (205.46,19.61) and (210.1,19.68) .. (212.98,22.59) .. controls (215.85,25.5) and (215.87,30.15) .. (213.01,32.97) .. controls (210.16,35.79) and (205.51,35.72) .. (202.63,32.8) .. controls (199.76,29.89) and (199.74,25.25) .. (202.6,22.43) -- cycle ;
%Shape: Ellipse [id:dp5733031479593196] 
\draw  [draw opacity=0][fill={rgb, 255:red, 15; green, 118; blue, 237 }  ,fill opacity=1 ] (255.6,22.43) .. controls (258.46,19.61) and (263.1,19.68) .. (265.98,22.59) .. controls (268.85,25.5) and (268.87,30.15) .. (266.01,32.97) .. controls (263.16,35.79) and (258.51,35.72) .. (255.63,32.8) .. controls (252.76,29.89) and (252.74,25.25) .. (255.6,22.43) -- cycle ;
%Shape: Ellipse [id:dp23658157144782477] 
\draw  [draw opacity=0][fill={rgb, 255:red, 15; green, 118; blue, 237 }  ,fill opacity=1 ] (310.6,22.43) .. controls (313.46,19.61) and (318.1,19.68) .. (320.98,22.59) .. controls (323.85,25.5) and (323.87,30.15) .. (321.01,32.97) .. controls (318.16,35.79) and (313.51,35.72) .. (310.63,32.8) .. controls (307.76,29.89) and (307.74,25.25) .. (310.6,22.43) -- cycle ;
%Shape: Ellipse [id:dp6142401450132564] 
\draw  [draw opacity=0][fill={rgb, 255:red, 15; green, 118; blue, 237 }  ,fill opacity=1 ] (202.6,94.43) .. controls (205.46,91.61) and (210.1,91.68) .. (212.98,94.59) .. controls (215.85,97.5) and (215.87,102.15) .. (213.01,104.97) .. controls (210.16,107.79) and (205.51,107.72) .. (202.63,104.8) .. controls (199.76,101.89) and (199.74,97.25) .. (202.6,94.43) -- cycle ;
%Shape: Ellipse [id:dp5054466614963726] 
\draw  [draw opacity=0][fill={rgb, 255:red, 15; green, 118; blue, 237 }  ,fill opacity=1 ] (255.6,94.43) .. controls (258.46,91.61) and (263.1,91.68) .. (265.98,94.59) .. controls (268.85,97.5) and (268.87,102.15) .. (266.01,104.97) .. controls (263.16,107.79) and (258.51,107.72) .. (255.63,104.8) .. controls (252.76,101.89) and (252.74,97.25) .. (255.6,94.43) -- cycle ;
%Shape: Ellipse [id:dp8994787904049966] 
\draw  [draw opacity=0][fill={rgb, 255:red, 15; green, 118; blue, 237 }  ,fill opacity=1 ] (310.6,94.43) .. controls (313.46,91.61) and (318.1,91.68) .. (320.98,94.59) .. controls (323.85,97.5) and (323.87,102.15) .. (321.01,104.97) .. controls (318.16,107.79) and (313.51,107.72) .. (310.63,104.8) .. controls (307.76,101.89) and (307.74,97.25) .. (310.6,94.43) -- cycle ;
%Straight Lines [id:da35684933345154146] 
\draw [color={rgb, 255:red, 16; green, 121; blue, 243 }  ,draw opacity=1 ][fill={rgb, 255:red, 0; green, 0; blue, 0 }  ,fill opacity=1 ][line width=1.5]    (207.81,27.7) -- (207.81,99.7) ;
%Straight Lines [id:da9025525649715185] 
\draw [color={rgb, 255:red, 16; green, 121; blue, 243 }  ,draw opacity=1 ][fill={rgb, 255:red, 0; green, 0; blue, 0 }  ,fill opacity=1 ][line width=1.5]    (207.81,27.7) -- (260.81,99.7) ;
%Straight Lines [id:da6066866708627501] 
\draw [color={rgb, 255:red, 16; green, 121; blue, 243 }  ,draw opacity=1 ][fill={rgb, 255:red, 0; green, 0; blue, 0 }  ,fill opacity=1 ][line width=1.5]    (207.81,27.7) -- (315.81,99.7) ;
%Straight Lines [id:da06743673012263462] 
\draw [color={rgb, 255:red, 16; green, 121; blue, 243 }  ,draw opacity=1 ][fill={rgb, 255:red, 0; green, 0; blue, 0 }  ,fill opacity=1 ][line width=1.5]    (260.81,27.7) -- (207.81,99.7) ;
%Straight Lines [id:da650936193176473] 
\draw [color={rgb, 255:red, 16; green, 121; blue, 243 }  ,draw opacity=1 ][fill={rgb, 255:red, 0; green, 0; blue, 0 }  ,fill opacity=1 ][line width=1.5]    (260.81,27.7) -- (260.81,99.7) ;
%Straight Lines [id:da6056347468202512] 
\draw [color={rgb, 255:red, 16; green, 121; blue, 243 }  ,draw opacity=1 ][fill={rgb, 255:red, 0; green, 0; blue, 0 }  ,fill opacity=1 ][line width=1.5]    (260.81,27.7) -- (315.81,99.7) ;
%Straight Lines [id:da3568672086228091] 
\draw [color={rgb, 255:red, 16; green, 121; blue, 243 }  ,draw opacity=1 ][fill={rgb, 255:red, 0; green, 0; blue, 0 }  ,fill opacity=1 ][line width=1.5]    (315.81,27.7) -- (207.81,99.7) ;
%Straight Lines [id:da8067833115504472] 
\draw [color={rgb, 255:red, 16; green, 121; blue, 243 }  ,draw opacity=1 ][fill={rgb, 255:red, 0; green, 0; blue, 0 }  ,fill opacity=1 ][line width=1.5]    (315.81,27.7) -- (260.81,99.7) ;
%Straight Lines [id:da7477661894800622] 
\draw [color={rgb, 255:red, 16; green, 121; blue, 243 }  ,draw opacity=1 ][fill={rgb, 255:red, 0; green, 0; blue, 0 }  ,fill opacity=1 ][line width=1.5]    (315.81,27.7) -- (315.81,99.7) ;

% Text Node
\draw (64.39,148.44) node [anchor=north west][inner sep=0.75pt]  [font=\small] [align=left] {(a)};
% Text Node
\draw (26,2.4) node [anchor=north west][inner sep=0.75pt]    {$a$};
% Text Node
\draw (9,90.4) node [anchor=north west][inner sep=0.75pt]    {$b$};
% Text Node
\draw (112,90.4) node [anchor=north west][inner sep=0.75pt]    {$c$};
% Text Node
\draw (110,19.4) node [anchor=north west][inner sep=0.75pt]    {$d$};
% Text Node
\draw (64,39.4) node [anchor=north west][inner sep=0.75pt]    {$e$};
% Text Node
\draw (146,52.4) node [anchor=north west][inner sep=0.75pt]    {$f$};
% Text Node
\draw (205,4.4) node [anchor=north west][inner sep=0.75pt]    {$a$};
% Text Node
\draw (206.63,110.2) node [anchor=north west][inner sep=0.75pt]    {$b$};
% Text Node
\draw (258,3.4) node [anchor=north west][inner sep=0.75pt]    {$c$};
% Text Node
\draw (314.63,110.2) node [anchor=north west][inner sep=0.75pt]    {$d$};
% Text Node
\draw (259.63,111.2) node [anchor=north west][inner sep=0.75pt]    {$e$};
% Text Node
\draw (312,5.4) node [anchor=north west][inner sep=0.75pt]    {$f$};
% Text Node
\draw (255.39,147.44) node [anchor=north west][inner sep=0.75pt]  [font=\small] [align=left] {(b)};


\end{tikzpicture}

\end{ftst}

%=====

\begin{ftst}{Grafos isomorfos}{Tipos de Grafos}
\begin{itemize}
    \item Dois grafos podem parecer diferentes em sua representação visual, mas ainda assim serem o mesmo grafo segundo a definição formal.\\
    \item O importante é saber distinguir entre dois grafos que têm estruturas fundamentais diferentes.
\end{itemize}

\vone
\centering


\tikzset{every picture/.style={line width=0.75pt}} %set default line width to 0.75pt        

\begin{tikzpicture}[x=0.75pt,y=0.75pt,yscale=-1,xscale=1]
%uncomment if require: \path (0,173); %set diagram left start at 0, and has height of 173

%Shape: Ellipse [id:dp7672867321663941] 
\draw  [draw opacity=0][fill={rgb, 255:red, 15; green, 118; blue, 237 }  ,fill opacity=1 ] (20.6,20.43) .. controls (23.46,17.61) and (28.1,17.68) .. (30.98,20.59) .. controls (33.85,23.5) and (33.87,28.15) .. (31.01,30.97) .. controls (28.16,33.79) and (23.51,33.72) .. (20.63,30.8) .. controls (17.76,27.89) and (17.74,23.25) .. (20.6,20.43) -- cycle ;
%Shape: Ellipse [id:dp7060944758385839] 
\draw  [draw opacity=0][fill={rgb, 255:red, 15; green, 118; blue, 237 }  ,fill opacity=1 ] (20.6,92.43) .. controls (23.46,89.61) and (28.1,89.68) .. (30.98,92.59) .. controls (33.85,95.5) and (33.87,100.15) .. (31.01,102.97) .. controls (28.16,105.79) and (23.51,105.72) .. (20.63,102.8) .. controls (17.76,99.89) and (17.74,95.25) .. (20.6,92.43) -- cycle ;
%Straight Lines [id:da6791312533711067] 
\draw [color={rgb, 255:red, 16; green, 121; blue, 243 }  ,draw opacity=1 ][fill={rgb, 255:red, 0; green, 0; blue, 0 }  ,fill opacity=1 ][line width=1.5]    (25.81,25.7) -- (25.81,97.7) ;
%Shape: Ellipse [id:dp7766028650957342] 
\draw  [draw opacity=0][fill={rgb, 255:red, 15; green, 118; blue, 237 }  ,fill opacity=1 ] (87.6,21.43) .. controls (90.46,18.61) and (95.1,18.68) .. (97.98,21.59) .. controls (100.85,24.5) and (100.87,29.15) .. (98.01,31.97) .. controls (95.16,34.79) and (90.51,34.72) .. (87.63,31.8) .. controls (84.76,28.89) and (84.74,24.25) .. (87.6,21.43) -- cycle ;
%Shape: Ellipse [id:dp13385211505551253] 
\draw  [draw opacity=0][fill={rgb, 255:red, 15; green, 118; blue, 237 }  ,fill opacity=1 ] (87.6,93.43) .. controls (90.46,90.61) and (95.1,90.68) .. (97.98,93.59) .. controls (100.85,96.5) and (100.87,101.15) .. (98.01,103.97) .. controls (95.16,106.79) and (90.51,106.72) .. (87.63,103.8) .. controls (84.76,100.89) and (84.74,96.25) .. (87.6,93.43) -- cycle ;
%Straight Lines [id:da7874419899288638] 
\draw [color={rgb, 255:red, 16; green, 121; blue, 243 }  ,draw opacity=1 ][fill={rgb, 255:red, 0; green, 0; blue, 0 }  ,fill opacity=1 ][line width=1.5]    (92.81,26.7) -- (92.81,98.7) ;
%Shape: Ellipse [id:dp3478384831309793] 
\draw  [draw opacity=0][fill={rgb, 255:red, 15; green, 118; blue, 237 }  ,fill opacity=1 ] (188.85,24.75) .. controls (189.5,20.79) and (193.29,18.11) .. (197.33,18.76) .. controls (201.37,19.42) and (204.12,23.16) .. (203.48,27.12) .. controls (202.84,31.08) and (199.04,33.76) .. (195,33.11) .. controls (190.97,32.45) and (188.21,28.71) .. (188.85,24.75) -- cycle ;
%Shape: Ellipse [id:dp6943800404481995] 
\draw  [draw opacity=0][fill={rgb, 255:red, 15; green, 118; blue, 237 }  ,fill opacity=1 ] (273.33,92.89) .. controls (273.97,88.92) and (277.77,86.24) .. (281.8,86.9) .. controls (285.84,87.55) and (288.6,91.3) .. (287.95,95.26) .. controls (287.31,99.22) and (283.52,101.9) .. (279.48,101.25) .. controls (275.44,100.59) and (272.69,96.85) .. (273.33,92.89) -- cycle ;
%Straight Lines [id:da30491395610387206] 
\draw [color={rgb, 255:red, 16; green, 121; blue, 243 }  ,draw opacity=1 ][fill={rgb, 255:red, 0; green, 0; blue, 0 }  ,fill opacity=1 ][line width=1.5]    (196.17,25.93) -- (280.64,94.07) ;
%Shape: Ellipse [id:dp22454624869229134] 
\draw  [draw opacity=0][fill={rgb, 255:red, 15; green, 118; blue, 237 }  ,fill opacity=1 ] (283.78,18.95) .. controls (287.75,19.58) and (290.44,23.37) .. (289.79,27.41) .. controls (289.15,31.45) and (285.42,34.21) .. (281.45,33.58) .. controls (277.49,32.95) and (274.8,29.16) .. (275.44,25.12) .. controls (276.08,21.08) and (279.82,18.32) .. (283.78,18.95) -- cycle ;
%Shape: Ellipse [id:dp2585636244259957] 
\draw  [draw opacity=0][fill={rgb, 255:red, 15; green, 118; blue, 237 }  ,fill opacity=1 ] (197.77,87.59) .. controls (201.73,88.22) and (204.43,92.01) .. (203.78,96.05) .. controls (203.14,100.09) and (199.41,102.86) .. (195.44,102.23) .. controls (191.48,101.6) and (188.78,97.81) .. (189.43,93.77) .. controls (190.07,89.73) and (193.8,86.96) .. (197.77,87.59) -- cycle ;
%Straight Lines [id:da09856693560407126] 
\draw [color={rgb, 255:red, 16; green, 121; blue, 243 }  ,draw opacity=1 ][fill={rgb, 255:red, 0; green, 0; blue, 0 }  ,fill opacity=1 ][line width=1.5]    (282.62,26.26) -- (196.61,94.91) ;

% Text Node
\draw (23,2.4) node [anchor=north west][inner sep=0.75pt]    {$v_{3}$};
% Text Node
\draw (24.63,108.2) node [anchor=north west][inner sep=0.75pt]    {$v_{1}$};
% Text Node
\draw (35,55.4) node [anchor=north west][inner sep=0.75pt]    {$a_{1}$};
% Text Node
\draw (90,3.4) node [anchor=north west][inner sep=0.75pt]    {$v_{2}$};
% Text Node
\draw (91.63,109.2) node [anchor=north west][inner sep=0.75pt]    {$v_{4}$};
% Text Node
\draw (102,56.4) node [anchor=north west][inner sep=0.75pt]    {$a_{2}$};
% Text Node
\draw (293.55,10.54) node [anchor=north west][inner sep=0.75pt]  [rotate=-358.99]  {$v_{3}$};
% Text Node
\draw (172.63,87.65) node [anchor=north west][inner sep=0.75pt]  [rotate=-1.25]  {$v_{1}$};
% Text Node
\draw (210.87,81.28) node [anchor=north west][inner sep=0.75pt]  [rotate=-0.73]  {$a_{1}$};
% Text Node
\draw (172.63,9.07) node [anchor=north west][inner sep=0.75pt]  [rotate=-0.52]  {$v_{2}$};
% Text Node
\draw (284.63,99.48) node [anchor=north west][inner sep=0.75pt]  [rotate=-0.16]  {$v_{4}$};
% Text Node
\draw (212.96,24.01) node [anchor=north west][inner sep=0.75pt]  [rotate=-358.65]  {$a_{2}$};
% Text Node
\draw (42.39,143.44) node [anchor=north west][inner sep=0.75pt]  [font=\small] [align=left] {Grafo G};
% Text Node
\draw (220.39,144.44) node [anchor=north west][inner sep=0.75pt]  [font=\small] [align=left] {Grafo H};


\end{tikzpicture}


\end{ftst}

%=====

\begin{ftst}{Grafos isomorfos}{Tipos de Grafos}

Para mostrar que duas estruturas são isomorfas, precisamos:
\vone
\begin{itemize}
    \item Obter uma bijeção entre os elementos das duas estruturas.
    \item Mostrar que as suas propriedades fundamentais são preservadas.
\end{itemize}

\vone

No caso dos grafos, os elementos são vértices e arestas, e a propriedade fundamental é “quais arestas conectam quais vértices” .
\end{ftst}


%=====

\begin{ftst}{Grafos isomorfos}{Tipos de Grafos}

Os grafos $G$ e $H$ a seguir são, essencialmente, o mesmo grafo, considerando as bijeções:
\vone



\tikzset{every picture/.style={line width=0.75pt}} %set default line width to 0.75pt        

\begin{tikzpicture}[x=0.75pt,y=0.75pt,yscale=-1,xscale=1]
%uncomment if require: \path (0,187); %set diagram left start at 0, and has height of 187

%Curve Lines [id:da13535893571863644] 
\draw [color={rgb, 255:red, 15; green, 118; blue, 237 }  ,draw opacity=1 ][line width=1.5]    (221.9,92.91) .. controls (267.74,83.37) and (275.04,43.91) .. (242.59,28.44) .. controls (210.14,12.96) and (161.46,14.51) .. (133.06,33.08) .. controls (104.67,51.65) and (109.53,79.5) .. (163.49,93.38) ;
%Shape: Ellipse [id:dp4051450804535437] 
\draw  [draw opacity=0][fill={rgb, 255:red, 15; green, 118; blue, 237 }  ,fill opacity=1 ] (14.2,23.2) .. controls (16.51,21.01) and (20.28,21.07) .. (22.62,23.32) .. controls (24.95,25.58) and (24.96,29.17) .. (22.64,31.35) .. controls (20.33,33.54) and (16.56,33.48) .. (14.23,31.23) .. controls (11.89,28.97) and (11.88,25.38) .. (14.2,23.2) -- cycle ;
%Shape: Ellipse [id:dp9979632274073846] 
\draw  [draw opacity=0][fill={rgb, 255:red, 15; green, 118; blue, 237 }  ,fill opacity=1 ] (14.2,78.9) .. controls (16.51,76.72) and (20.28,76.78) .. (22.62,79.03) .. controls (24.95,81.28) and (24.96,84.88) .. (22.64,87.06) .. controls (20.33,89.24) and (16.56,89.19) .. (14.23,86.93) .. controls (11.89,84.68) and (11.88,81.09) .. (14.2,78.9) -- cycle ;
%Straight Lines [id:da7110634067887398] 
\draw [color={rgb, 255:red, 16; green, 121; blue, 243 }  ,draw opacity=1 ][fill={rgb, 255:red, 0; green, 0; blue, 0 }  ,fill opacity=1 ][line width=1.5]    (18.42,27.27) -- (18.42,82.98) ;
%Shape: Ellipse [id:dp37656452184529754] 
\draw  [draw opacity=0][fill={rgb, 255:red, 15; green, 118; blue, 237 }  ,fill opacity=1 ] (68.56,23.97) .. controls (70.87,21.79) and (74.64,21.84) .. (76.98,24.1) .. controls (79.31,26.35) and (79.32,29.94) .. (77,32.13) .. controls (74.69,34.31) and (70.92,34.25) .. (68.58,32) .. controls (66.25,29.75) and (66.24,26.15) .. (68.56,23.97) -- cycle ;
%Shape: Ellipse [id:dp8237350008256439] 
\draw  [draw opacity=0][fill={rgb, 255:red, 15; green, 118; blue, 237 }  ,fill opacity=1 ] (68.56,79.68) .. controls (70.87,77.49) and (74.64,77.55) .. (76.98,79.8) .. controls (79.31,82.06) and (79.32,85.65) .. (77,87.83) .. controls (74.69,90.02) and (70.92,89.96) .. (68.58,87.71) .. controls (66.25,85.45) and (66.24,81.86) .. (68.56,79.68) -- cycle ;
%Straight Lines [id:da9589476920573761] 
\draw [color={rgb, 255:red, 16; green, 121; blue, 243 }  ,draw opacity=1 ][fill={rgb, 255:red, 0; green, 0; blue, 0 }  ,fill opacity=1 ][line width=1.5]    (72.78,28.05) -- (72.78,83.76) ;
%Shape: Ellipse [id:dp7592445989982248] 
\draw  [draw opacity=0][fill={rgb, 255:red, 15; green, 118; blue, 237 }  ,fill opacity=1 ] (158.43,61.85) .. controls (156.13,59.66) and (156.15,56.07) .. (158.49,53.82) .. controls (160.84,51.58) and (164.61,51.54) .. (166.91,53.73) .. controls (169.22,55.92) and (169.19,59.51) .. (166.85,61.76) .. controls (164.51,64) and (160.74,64.04) .. (158.43,61.85) -- cycle ;
%Shape: Ellipse [id:dp46395815979252064] 
\draw  [draw opacity=0][fill={rgb, 255:red, 15; green, 118; blue, 237 }  ,fill opacity=1 ] (216.85,61.38) .. controls (214.54,59.19) and (214.57,55.59) .. (216.91,53.35) .. controls (219.25,51.1) and (223.02,51.06) .. (225.33,53.25) .. controls (227.64,55.44) and (227.61,59.04) .. (225.27,61.28) .. controls (222.92,63.53) and (219.15,63.57) .. (216.85,61.38) -- cycle ;
%Straight Lines [id:da10247453238790105] 
\draw [color={rgb, 255:red, 16; green, 121; blue, 243 }  ,draw opacity=1 ][fill={rgb, 255:red, 0; green, 0; blue, 0 }  ,fill opacity=1 ][line width=1.5]    (162.67,57.79) -- (221.09,57.31) ;
%Shape: Ellipse [id:dp3514904288568135] 
\draw  [draw opacity=0][fill={rgb, 255:red, 15; green, 118; blue, 237 }  ,fill opacity=1 ] (159.24,95.9) .. controls (156.94,93.71) and (156.96,90.11) .. (159.31,87.87) .. controls (161.65,85.62) and (165.42,85.58) .. (167.73,87.77) .. controls (170.03,89.96) and (170.01,93.56) .. (167.66,95.8) .. controls (165.32,98.05) and (161.55,98.09) .. (159.24,95.9) -- cycle ;
%Shape: Ellipse [id:dp7592627436182027] 
\draw  [draw opacity=0][fill={rgb, 255:red, 15; green, 118; blue, 237 }  ,fill opacity=1 ] (217.66,95.42) .. controls (215.35,93.23) and (215.38,89.63) .. (217.72,87.39) .. controls (220.06,85.15) and (223.83,85.1) .. (226.14,87.3) .. controls (228.45,89.49) and (228.42,93.08) .. (226.08,95.33) .. controls (223.73,97.57) and (219.97,97.61) .. (217.66,95.42) -- cycle ;

% Text Node
\draw (14.82,7.76) node [anchor=north west][inner sep=0.75pt]    {$v_{3}$};
% Text Node
\draw (16.15,89.62) node [anchor=north west][inner sep=0.75pt]    {$v_{1}$};
% Text Node
\draw (24.65,48.76) node [anchor=north west][inner sep=0.75pt]    {$a_{1}$};
% Text Node
\draw (69.18,8.53) node [anchor=north west][inner sep=0.75pt]    {$v_{2}$};
% Text Node
\draw (70.51,90.39) node [anchor=north west][inner sep=0.75pt]    {$v_{4}$};
% Text Node
\draw (79.01,49.54) node [anchor=north west][inner sep=0.75pt]    {$a_{2}$};
% Text Node
\draw (27.53,119.56) node [anchor=north west][inner sep=0.75pt]  [font=\small] [align=left] {Grafo G};
% Text Node
\draw (140.6,48.8) node [anchor=north west][inner sep=0.75pt]  [rotate=-0.23]  {$v_{a}$};
% Text Node
\draw (229.62,48.34) node [anchor=north west][inner sep=0.75pt]  [rotate=-0.45]  {$v_{b}$};
% Text Node
\draw (185.57,38.85) node [anchor=north west][inner sep=0.75pt]  [rotate=-359.74]  {$e_{1}$};
% Text Node
\draw (137.7,7.9) node [anchor=north west][inner sep=0.75pt]  [rotate=-359.74]  {$e_{2}$};
% Text Node
\draw (156.83,96.77) node [anchor=north west][inner sep=0.75pt]  [rotate=-0.23]  {$v_{c}$};
% Text Node
\draw (215.52,96.56) node [anchor=north west][inner sep=0.75pt]  [rotate=-0.23]  {$v_{d}$};
% Text Node
\draw (169.51,120.76) node [anchor=north west][inner sep=0.75pt]  [font=\small] [align=left] {Grafo H};
% Text Node
\draw (285.35,20.41) node [anchor=north west][inner sep=0.75pt]  [font=\small] [align=left] {Bijeção entre os vértices: $\displaystyle s$};
% Text Node
\draw (286.41,96.79) node [anchor=north west][inner sep=0.75pt]  [font=\small] [align=left] {Bijeção entre as arestas: $\displaystyle t$};
% Text Node
\draw (305.74,44.17) node [anchor=north west][inner sep=0.75pt]    {$v_{1} \ \Leftrightarrow \ v_{a}$};
% Text Node
\draw (304.8,69.33) node [anchor=north west][inner sep=0.75pt]    {$v_{3} \ \Leftrightarrow \ v_{b}$};
% Text Node
\draw (406.88,44.17) node [anchor=north west][inner sep=0.75pt]    {$v_{2} \ \Leftrightarrow \ v_{c}$};
% Text Node
\draw (405.91,69.33) node [anchor=north west][inner sep=0.75pt]    {$v_{4} \ \Leftrightarrow \ v_{d}$};
% Text Node
\draw (304.9,120.55) node [anchor=north west][inner sep=0.75pt]    {$a_{1} \ \Leftrightarrow \ e_{1}$};
% Text Node
\draw (303.96,145.71) node [anchor=north west][inner sep=0.75pt]    {$a_{2} \ \Leftrightarrow \ e_{2}$};


\end{tikzpicture}

\end{ftst}


%=====

\begin{ftst}{Grafos isomorfos}{Tipos de Grafos}

Exemplos de condições nas quais dois grafos não são isomorfos:
\vone
\begin{itemize}
    \item Um grafo tem mais vértices que o outro;
    \item Um grafo tem mais arestas que o outro;
    \item Um grafo tem arcos paralelos, um laço ou um ciclo, e o outro não;
    \item Um grafo tem um vértice de grau k e o outro não;
    \item Um grafo é conexo e outro não, etc. 
\end{itemize}

\end{ftst}


%=====

\begin{ftst}{Exercícios}{Tipos de Grafos}

Os grafos $G$ e $H$ a seguir são isomorfos? Se sim, apresente as bijeções que estabelecem o isomorfismo.

\vone
\centering


\tikzset{every picture/.style={line width=0.75pt}} %set default line width to 0.75pt        

\begin{tikzpicture}[x=0.75pt,y=0.75pt,yscale=-1,xscale=1]
%uncomment if require: \path (0,197); %set diagram left start at 0, and has height of 197

%Shape: Circle [id:dp09163317413291105] 
\draw  [draw opacity=0][fill={rgb, 255:red, 15; green, 118; blue, 237 }  ,fill opacity=1 ] (75,29.8) .. controls (75,26.6) and (77.6,24) .. (80.8,24) .. controls (84.01,24) and (86.61,26.6) .. (86.61,29.8) .. controls (86.61,33.01) and (84.01,35.61) .. (80.8,35.61) .. controls (77.6,35.61) and (75,33.01) .. (75,29.8) -- cycle ;
%Shape: Circle [id:dp8468757243290981] 
\draw  [draw opacity=0][fill={rgb, 255:red, 15; green, 118; blue, 237 }  ,fill opacity=1 ] (122,29.8) .. controls (122,26.6) and (124.6,24) .. (127.8,24) .. controls (131.01,24) and (133.61,26.6) .. (133.61,29.8) .. controls (133.61,33.01) and (131.01,35.61) .. (127.8,35.61) .. controls (124.6,35.61) and (122,33.01) .. (122,29.8) -- cycle ;
%Shape: Circle [id:dp5182730656612713] 
\draw  [draw opacity=0][fill={rgb, 255:red, 15; green, 118; blue, 237 }  ,fill opacity=1 ] (75,126.8) .. controls (75,123.6) and (77.6,121) .. (80.8,121) .. controls (84.01,121) and (86.61,123.6) .. (86.61,126.8) .. controls (86.61,130.01) and (84.01,132.61) .. (80.8,132.61) .. controls (77.6,132.61) and (75,130.01) .. (75,126.8) -- cycle ;
%Shape: Circle [id:dp2348776010167375] 
\draw  [draw opacity=0][fill={rgb, 255:red, 15; green, 118; blue, 237 }  ,fill opacity=1 ] (122,126.8) .. controls (122,123.6) and (124.6,121) .. (127.8,121) .. controls (131.01,121) and (133.61,123.6) .. (133.61,126.8) .. controls (133.61,130.01) and (131.01,132.61) .. (127.8,132.61) .. controls (124.6,132.61) and (122,130.01) .. (122,126.8) -- cycle ;
%Shape: Circle [id:dp8699381344171488] 
\draw  [draw opacity=0][fill={rgb, 255:red, 15; green, 118; blue, 237 }  ,fill opacity=1 ] (27,77.8) .. controls (27,74.6) and (29.6,72) .. (32.8,72) .. controls (36.01,72) and (38.61,74.6) .. (38.61,77.8) .. controls (38.61,81.01) and (36.01,83.61) .. (32.8,83.61) .. controls (29.6,83.61) and (27,81.01) .. (27,77.8) -- cycle ;
%Shape: Circle [id:dp6385775932083995] 
\draw  [draw opacity=0][fill={rgb, 255:red, 15; green, 118; blue, 237 }  ,fill opacity=1 ] (171,77.8) .. controls (171,74.6) and (173.6,72) .. (176.8,72) .. controls (180.01,72) and (182.61,74.6) .. (182.61,77.8) .. controls (182.61,81.01) and (180.01,83.61) .. (176.8,83.61) .. controls (173.6,83.61) and (171,81.01) .. (171,77.8) -- cycle ;
%Shape: Circle [id:dp8611713488366561] 
\draw  [draw opacity=0][fill={rgb, 255:red, 15; green, 118; blue, 237 }  ,fill opacity=1 ] (99,77.8) .. controls (99,74.6) and (101.6,72) .. (104.8,72) .. controls (108.01,72) and (110.61,74.6) .. (110.61,77.8) .. controls (110.61,81.01) and (108.01,83.61) .. (104.8,83.61) .. controls (101.6,83.61) and (99,81.01) .. (99,77.8) -- cycle ;
%Straight Lines [id:da33404860781636514] 
\draw [color={rgb, 255:red, 16; green, 121; blue, 243 }  ,draw opacity=1 ][fill={rgb, 255:red, 0; green, 0; blue, 0 }  ,fill opacity=1 ][line width=1.5]    (80.8,29.8) -- (32.8,77.8) ;
%Straight Lines [id:da29430790395343753] 
\draw [color={rgb, 255:red, 16; green, 121; blue, 243 }  ,draw opacity=1 ][fill={rgb, 255:red, 0; green, 0; blue, 0 }  ,fill opacity=1 ][line width=1.5]    (32.8,77.8) -- (80.8,126.8) ;
%Straight Lines [id:da2053532342114761] 
\draw [color={rgb, 255:red, 16; green, 121; blue, 243 }  ,draw opacity=1 ][fill={rgb, 255:red, 0; green, 0; blue, 0 }  ,fill opacity=1 ][line width=1.5]    (127.8,126.8) -- (80.8,126.8) ;
%Straight Lines [id:da31411890227171546] 
\draw [color={rgb, 255:red, 16; green, 121; blue, 243 }  ,draw opacity=1 ][fill={rgb, 255:red, 0; green, 0; blue, 0 }  ,fill opacity=1 ][line width=1.5]    (127.8,29.8) -- (80.8,29.8) ;
%Straight Lines [id:da0953839058044399] 
\draw [color={rgb, 255:red, 16; green, 121; blue, 243 }  ,draw opacity=1 ][fill={rgb, 255:red, 0; green, 0; blue, 0 }  ,fill opacity=1 ][line width=1.5]    (127.8,29.8) -- (176.8,77.8) ;
%Straight Lines [id:da8305904330512186] 
\draw [color={rgb, 255:red, 16; green, 121; blue, 243 }  ,draw opacity=1 ][fill={rgb, 255:red, 0; green, 0; blue, 0 }  ,fill opacity=1 ][line width=1.5]    (176.8,77.8) -- (127.8,126.8) ;
%Straight Lines [id:da21762196114326238] 
\draw [color={rgb, 255:red, 16; green, 121; blue, 243 }  ,draw opacity=1 ][fill={rgb, 255:red, 0; green, 0; blue, 0 }  ,fill opacity=1 ][line width=1.5]    (80.8,29.8) -- (104.8,77.8) ;
%Straight Lines [id:da9173257905513856] 
\draw [color={rgb, 255:red, 16; green, 121; blue, 243 }  ,draw opacity=1 ][fill={rgb, 255:red, 0; green, 0; blue, 0 }  ,fill opacity=1 ][line width=1.5]    (104.8,77.8) -- (127.8,126.8) ;
%Straight Lines [id:da9944257200619819] 
\draw [color={rgb, 255:red, 16; green, 121; blue, 243 }  ,draw opacity=1 ][fill={rgb, 255:red, 0; green, 0; blue, 0 }  ,fill opacity=1 ][line width=1.5]    (176.8,77.8) -- (104.8,77.8) ;
%Shape: Circle [id:dp5818714406656633] 
\draw  [draw opacity=0][fill={rgb, 255:red, 15; green, 118; blue, 237 }  ,fill opacity=1 ] (284,29.8) .. controls (284,26.6) and (286.6,24) .. (289.8,24) .. controls (293.01,24) and (295.61,26.6) .. (295.61,29.8) .. controls (295.61,33.01) and (293.01,35.61) .. (289.8,35.61) .. controls (286.6,35.61) and (284,33.01) .. (284,29.8) -- cycle ;
%Shape: Circle [id:dp449760456790973] 
\draw  [draw opacity=0][fill={rgb, 255:red, 15; green, 118; blue, 237 }  ,fill opacity=1 ] (331,29.8) .. controls (331,26.6) and (333.6,24) .. (336.8,24) .. controls (340.01,24) and (342.61,26.6) .. (342.61,29.8) .. controls (342.61,33.01) and (340.01,35.61) .. (336.8,35.61) .. controls (333.6,35.61) and (331,33.01) .. (331,29.8) -- cycle ;
%Shape: Circle [id:dp39197426344427333] 
\draw  [draw opacity=0][fill={rgb, 255:red, 15; green, 118; blue, 237 }  ,fill opacity=1 ] (380,29.8) .. controls (380,26.6) and (382.6,24) .. (385.8,24) .. controls (389.01,24) and (391.61,26.6) .. (391.61,29.8) .. controls (391.61,33.01) and (389.01,35.61) .. (385.8,35.61) .. controls (382.6,35.61) and (380,33.01) .. (380,29.8) -- cycle ;
%Shape: Circle [id:dp8872822023982192] 
\draw  [draw opacity=0][fill={rgb, 255:red, 15; green, 118; blue, 237 }  ,fill opacity=1 ] (283,125.8) .. controls (283,122.6) and (285.6,120) .. (288.8,120) .. controls (292.01,120) and (294.61,122.6) .. (294.61,125.8) .. controls (294.61,129.01) and (292.01,131.61) .. (288.8,131.61) .. controls (285.6,131.61) and (283,129.01) .. (283,125.8) -- cycle ;
%Shape: Circle [id:dp8111175509068969] 
\draw  [draw opacity=0][fill={rgb, 255:red, 15; green, 118; blue, 237 }  ,fill opacity=1 ] (330,125.8) .. controls (330,122.6) and (332.6,120) .. (335.8,120) .. controls (339.01,120) and (341.61,122.6) .. (341.61,125.8) .. controls (341.61,129.01) and (339.01,131.61) .. (335.8,131.61) .. controls (332.6,131.61) and (330,129.01) .. (330,125.8) -- cycle ;
%Shape: Circle [id:dp5427454292962486] 
\draw  [draw opacity=0][fill={rgb, 255:red, 15; green, 118; blue, 237 }  ,fill opacity=1 ] (379,125.8) .. controls (379,122.6) and (381.6,120) .. (384.8,120) .. controls (388.01,120) and (390.61,122.6) .. (390.61,125.8) .. controls (390.61,129.01) and (388.01,131.61) .. (384.8,131.61) .. controls (381.6,131.61) and (379,129.01) .. (379,125.8) -- cycle ;
%Shape: Circle [id:dp4798004440975898] 
\draw  [draw opacity=0][fill={rgb, 255:red, 15; green, 118; blue, 237 }  ,fill opacity=1 ] (331,77.8) .. controls (331,74.6) and (333.6,72) .. (336.8,72) .. controls (340.01,72) and (342.61,74.6) .. (342.61,77.8) .. controls (342.61,81.01) and (340.01,83.61) .. (336.8,83.61) .. controls (333.6,83.61) and (331,81.01) .. (331,77.8) -- cycle ;
%Straight Lines [id:da5065872337975126] 
\draw [color={rgb, 255:red, 16; green, 121; blue, 243 }  ,draw opacity=1 ][fill={rgb, 255:red, 0; green, 0; blue, 0 }  ,fill opacity=1 ][line width=1.5]    (289.8,29.8) -- (288.8,125.8) ;
%Straight Lines [id:da5436441309271014] 
\draw [color={rgb, 255:red, 16; green, 121; blue, 243 }  ,draw opacity=1 ][fill={rgb, 255:red, 0; green, 0; blue, 0 }  ,fill opacity=1 ][line width=1.5]    (385.8,29.8) -- (384.8,125.8) ;
%Straight Lines [id:da5936877200307056] 
\draw [color={rgb, 255:red, 16; green, 121; blue, 243 }  ,draw opacity=1 ][fill={rgb, 255:red, 0; green, 0; blue, 0 }  ,fill opacity=1 ][line width=1.5]    (385.8,29.8) -- (289.8,29.8) ;
%Straight Lines [id:da4072879069803035] 
\draw [color={rgb, 255:red, 16; green, 121; blue, 243 }  ,draw opacity=1 ][fill={rgb, 255:red, 0; green, 0; blue, 0 }  ,fill opacity=1 ][line width=1.5]    (384.8,125.8) -- (288.8,125.8) ;
%Straight Lines [id:da6010724750209526] 
\draw [color={rgb, 255:red, 16; green, 121; blue, 243 }  ,draw opacity=1 ][fill={rgb, 255:red, 0; green, 0; blue, 0 }  ,fill opacity=1 ][line width=1.5]    (289.8,29.8) -- (335.8,125.8) ;
%Straight Lines [id:da8630680139814437] 
\draw [color={rgb, 255:red, 16; green, 121; blue, 243 }  ,draw opacity=1 ][fill={rgb, 255:red, 0; green, 0; blue, 0 }  ,fill opacity=1 ][line width=1.5]    (336.8,77.8) -- (288.8,125.8) ;
%Straight Lines [id:da8237346631176974] 
\draw [color={rgb, 255:red, 16; green, 121; blue, 243 }  ,draw opacity=1 ][fill={rgb, 255:red, 0; green, 0; blue, 0 }  ,fill opacity=1 ][line width=1.5]    (336.8,77.8) -- (384.8,125.8) ;

% Text Node
\draw (70,8.4) node [anchor=north west][inner sep=0.75pt]  [color={rgb, 255:red, 208; green, 2; blue, 27 }  ,opacity=1 ]  {$A$};
% Text Node
\draw (128,9.4) node [anchor=north west][inner sep=0.75pt]  [color={rgb, 255:red, 208; green, 2; blue, 27 }  ,opacity=1 ]  {$B$};
% Text Node
\draw (10,72.4) node [anchor=north west][inner sep=0.75pt]  [color={rgb, 255:red, 208; green, 2; blue, 27 }  ,opacity=1 ]  {$C$};
% Text Node
\draw (109,57.4) node [anchor=north west][inner sep=0.75pt]  [color={rgb, 255:red, 208; green, 2; blue, 27 }  ,opacity=1 ]  {$D$};
% Text Node
\draw (188,70.4) node [anchor=north west][inner sep=0.75pt]  [color={rgb, 255:red, 208; green, 2; blue, 27 }  ,opacity=1 ]  {$E$};
% Text Node
\draw (78,137.4) node [anchor=north west][inner sep=0.75pt]  [color={rgb, 255:red, 208; green, 2; blue, 27 }  ,opacity=1 ]  {$F$};
% Text Node
\draw (124,137.4) node [anchor=north west][inner sep=0.75pt]  [color={rgb, 255:red, 208; green, 2; blue, 27 }  ,opacity=1 ]  {$G$};
% Text Node
\draw (156.3,107.7) node [anchor=north west][inner sep=0.75pt]    {$a_{1}$};
% Text Node
\draw (119.01,90.54) node [anchor=north west][inner sep=0.75pt]    {$a_{2}$};
% Text Node
\draw (97.01,129.54) node [anchor=north west][inner sep=0.75pt]    {$a_{3}$};
% Text Node
\draw (42.01,103.54) node [anchor=north west][inner sep=0.75pt]    {$a_{4}$};
% Text Node
\draw (135.01,59.54) node [anchor=north west][inner sep=0.75pt]    {$a_{5}$};
% Text Node
\draw (153.01,34.54) node [anchor=north west][inner sep=0.75pt]    {$a_{6}$};
% Text Node
\draw (76.01,50.54) node [anchor=north west][inner sep=0.75pt]    {$a_{7}$};
% Text Node
\draw (98.01,11.54) node [anchor=north west][inner sep=0.75pt]    {$a_{8}$};
% Text Node
\draw (38.01,36.54) node [anchor=north west][inner sep=0.75pt]    {$a_{9}$};
% Text Node
\draw (286,7.4) node [anchor=north west][inner sep=0.75pt]  [color={rgb, 255:red, 208; green, 2; blue, 27 }  ,opacity=1 ]  {$1$};
% Text Node
\draw (333,5.4) node [anchor=north west][inner sep=0.75pt]  [color={rgb, 255:red, 208; green, 2; blue, 27 }  ,opacity=1 ]  {$2$};
% Text Node
\draw (382,6.4) node [anchor=north west][inner sep=0.75pt]  [color={rgb, 255:red, 208; green, 2; blue, 27 }  ,opacity=1 ]  {$3$};
% Text Node
\draw (334,56.4) node [anchor=north west][inner sep=0.75pt]  [color={rgb, 255:red, 208; green, 2; blue, 27 }  ,opacity=1 ]  {$4$};
% Text Node
\draw (285,138.4) node [anchor=north west][inner sep=0.75pt]  [color={rgb, 255:red, 208; green, 2; blue, 27 }  ,opacity=1 ]  {$5$};
% Text Node
\draw (332,138.4) node [anchor=north west][inner sep=0.75pt]  [color={rgb, 255:red, 208; green, 2; blue, 27 }  ,opacity=1 ]  {$6$};
% Text Node
\draw (381.8,138.01) node [anchor=north west][inner sep=0.75pt]  [color={rgb, 255:red, 208; green, 2; blue, 27 }  ,opacity=1 ]  {$7$};
% Text Node
\draw (271.3,64.7) node [anchor=north west][inner sep=0.75pt]    {$e_{1}$};
% Text Node
\draw (307.01,127.54) node [anchor=north west][inner sep=0.75pt]    {$e_{2}$};
% Text Node
\draw (354.01,128.54) node [anchor=north west][inner sep=0.75pt]    {$e_{3}$};
% Text Node
\draw (308.01,46.54) node [anchor=north west][inner sep=0.75pt]    {$e_{4}$};
% Text Node
\draw (301.01,82.54) node [anchor=north west][inner sep=0.75pt]    {$e_{5}$};
% Text Node
\draw (361.01,81.54) node [anchor=north west][inner sep=0.75pt]    {$e_{6}$};
% Text Node
\draw (391.01,65.54) node [anchor=north west][inner sep=0.75pt]    {$e_{7}$};
% Text Node
\draw (307.01,8.54) node [anchor=north west][inner sep=0.75pt]    {$e_{8}$};
% Text Node
\draw (356.01,8.54) node [anchor=north west][inner sep=0.75pt]    {$e_{9}$};
% Text Node
\draw (81.53,162.56) node [anchor=north west][inner sep=0.75pt]  [font=\small] [align=left] {Grafo G};
% Text Node
\draw (313.51,163.76) node [anchor=north west][inner sep=0.75pt]  [font=\small] [align=left] {Grafo H};


\end{tikzpicture}

\end{ftst}


%=====

\begin{ftst}{Exercícios}{Tipos de Grafos}

Os grafos $G$ e $H$ a seguir são isomorfos? Se sim, apresente as bijeções que estabelecem o isomorfismo.

\vone
\centering


\tikzset{every picture/.style={line width=0.75pt}} %set default line width to 0.75pt        

\begin{tikzpicture}[x=0.75pt,y=0.75pt,yscale=-1,xscale=1]
%uncomment if require: \path (0,238); %set diagram left start at 0, and has height of 238

%Shape: Ellipse [id:dp8399187232641918] 
\draw  [draw opacity=0][fill={rgb, 255:red, 15; green, 118; blue, 237 }  ,fill opacity=1 ] (143.36,36.95) .. controls (143.36,33.99) and (145.87,31.59) .. (148.96,31.59) .. controls (152.05,31.59) and (154.56,33.99) .. (154.56,36.95) .. controls (154.56,39.91) and (152.05,42.31) .. (148.96,42.31) .. controls (145.87,42.31) and (143.36,39.91) .. (143.36,36.95) -- cycle ;
%Shape: Ellipse [id:dp8527982659961908] 
\draw  [draw opacity=0][fill={rgb, 255:red, 15; green, 118; blue, 237 }  ,fill opacity=1 ] (88.38,169.99) .. controls (88.38,167.03) and (90.88,164.62) .. (93.98,164.62) .. controls (97.07,164.62) and (99.58,167.03) .. (99.58,169.99) .. controls (99.58,172.95) and (97.07,175.35) .. (93.98,175.35) .. controls (90.88,175.35) and (88.38,172.95) .. (88.38,169.99) -- cycle ;
%Shape: Ellipse [id:dp881675890625838] 
\draw  [draw opacity=0][fill={rgb, 255:red, 15; green, 118; blue, 237 }  ,fill opacity=1 ] (32.43,158.9) .. controls (32.43,155.94) and (34.94,153.54) .. (38.03,153.54) .. controls (41.12,153.54) and (43.63,155.94) .. (43.63,158.9) .. controls (43.63,161.86) and (41.12,164.26) .. (38.03,164.26) .. controls (34.94,164.26) and (32.43,161.86) .. (32.43,158.9) -- cycle ;
%Shape: Ellipse [id:dp9410795918956025] 
\draw  [draw opacity=0][fill={rgb, 255:red, 15; green, 118; blue, 237 }  ,fill opacity=1 ] (88.38,31.41) .. controls (88.38,28.44) and (90.88,26.04) .. (93.98,26.04) .. controls (97.07,26.04) and (99.58,28.44) .. (99.58,31.41) .. controls (99.58,34.37) and (97.07,36.77) .. (93.98,36.77) .. controls (90.88,36.77) and (88.38,34.37) .. (88.38,31.41) -- cycle ;
%Shape: Ellipse [id:dp03988826531518175] 
\draw  [draw opacity=0][fill={rgb, 255:red, 15; green, 118; blue, 237 }  ,fill opacity=1 ] (88.38,96.08) .. controls (88.38,93.12) and (90.88,90.71) .. (93.98,90.71) .. controls (97.07,90.71) and (99.58,93.12) .. (99.58,96.08) .. controls (99.58,99.04) and (97.07,101.44) .. (93.98,101.44) .. controls (90.88,101.44) and (88.38,99.04) .. (88.38,96.08) -- cycle ;
%Straight Lines [id:da43215622185878644] 
\draw [color={rgb, 255:red, 16; green, 121; blue, 243 }  ,draw opacity=1 ][fill={rgb, 255:red, 0; green, 0; blue, 0 }  ,fill opacity=1 ][line width=1.5]    (38.03,158.9) -- (148.96,36.95) ;
%Straight Lines [id:da800285207328959] 
\draw [color={rgb, 255:red, 16; green, 121; blue, 243 }  ,draw opacity=1 ][fill={rgb, 255:red, 0; green, 0; blue, 0 }  ,fill opacity=1 ][line width=1.5]    (93.98,32.79) -- (93.98,169.99) ;
%Curve Lines [id:da5214066995754527] 
\draw [color={rgb, 255:red, 16; green, 121; blue, 243 }  ,draw opacity=1 ][line width=1.5]    (93.98,32.79) .. controls (125.13,8.56) and (137.67,23.34) .. (148.96,36.95) ;
%Curve Lines [id:da5052043551765479] 
\draw [color={rgb, 255:red, 16; green, 121; blue, 243 }  ,draw opacity=1 ][line width=1.5]    (93.98,31.41) .. controls (108.74,55.67) and (132.85,48.28) .. (148.96,36.95) ;
%Curve Lines [id:da6690163807214244] 
\draw [color={rgb, 255:red, 16; green, 121; blue, 243 }  ,draw opacity=1 ][line width=1.5]    (38.03,158.9) .. controls (69.19,134.66) and (82.69,155.45) .. (93.98,169.06) ;
%Curve Lines [id:da05952989424985766] 
\draw [color={rgb, 255:red, 16; green, 121; blue, 243 }  ,draw opacity=1 ][line width=1.5]    (38.03,158.9) .. controls (52.79,183.17) and (77.87,180.4) .. (93.98,169.06) ;
%Shape: Ellipse [id:dp8144869141971312] 
\draw  [draw opacity=0][fill={rgb, 255:red, 15; green, 118; blue, 237 }  ,fill opacity=1 ] (285.15,32.33) .. controls (285.15,29.37) and (287.66,26.97) .. (290.75,26.97) .. controls (293.84,26.97) and (296.35,29.37) .. (296.35,32.33) .. controls (296.35,35.29) and (293.84,37.69) .. (290.75,37.69) .. controls (287.66,37.69) and (285.15,35.29) .. (285.15,32.33) -- cycle ;
%Shape: Ellipse [id:dp4147749868866344] 
\draw  [draw opacity=0][fill={rgb, 255:red, 15; green, 118; blue, 237 }  ,fill opacity=1 ] (285.15,98.85) .. controls (285.15,95.89) and (287.66,93.49) .. (290.75,93.49) .. controls (293.84,93.49) and (296.35,95.89) .. (296.35,98.85) .. controls (296.35,101.81) and (293.84,104.21) .. (290.75,104.21) .. controls (287.66,104.21) and (285.15,101.81) .. (285.15,98.85) -- cycle ;
%Shape: Ellipse [id:dp5616701178049268] 
\draw  [draw opacity=0][fill={rgb, 255:red, 15; green, 118; blue, 237 }  ,fill opacity=1 ] (285.15,161.67) .. controls (285.15,158.71) and (287.66,156.31) .. (290.75,156.31) .. controls (293.84,156.31) and (296.35,158.71) .. (296.35,161.67) .. controls (296.35,164.63) and (293.84,167.03) .. (290.75,167.03) .. controls (287.66,167.03) and (285.15,164.63) .. (285.15,161.67) -- cycle ;
%Shape: Ellipse [id:dp27707467924239926] 
\draw  [draw opacity=0][fill={rgb, 255:red, 15; green, 118; blue, 237 }  ,fill opacity=1 ] (326.63,98.85) .. controls (326.63,95.89) and (329.13,93.49) .. (332.23,93.49) .. controls (335.32,93.49) and (337.83,95.89) .. (337.83,98.85) .. controls (337.83,101.81) and (335.32,104.21) .. (332.23,104.21) .. controls (329.13,104.21) and (326.63,101.81) .. (326.63,98.85) -- cycle ;
%Shape: Ellipse [id:dp5877586830595478] 
\draw  [draw opacity=0][fill={rgb, 255:red, 15; green, 118; blue, 237 }  ,fill opacity=1 ] (244.64,97.92) .. controls (244.64,94.96) and (247.15,92.56) .. (250.24,92.56) .. controls (253.33,92.56) and (255.84,94.96) .. (255.84,97.92) .. controls (255.84,100.89) and (253.33,103.29) .. (250.24,103.29) .. controls (247.15,103.29) and (244.64,100.89) .. (244.64,97.92) -- cycle ;
%Straight Lines [id:da19146929368315546] 
\draw [color={rgb, 255:red, 16; green, 121; blue, 243 }  ,draw opacity=1 ][fill={rgb, 255:red, 0; green, 0; blue, 0 }  ,fill opacity=1 ][line width=1.5]    (290.75,161.67) -- (290.75,32.33) ;
%Straight Lines [id:da28566898674225527] 
\draw [color={rgb, 255:red, 16; green, 121; blue, 243 }  ,draw opacity=1 ][fill={rgb, 255:red, 0; green, 0; blue, 0 }  ,fill opacity=1 ][line width=1.5]    (250.24,97.92) -- (332.23,98.85) ;
%Shape: Ellipse [id:dp5216131110171041] 
\draw  [color={rgb, 255:red, 16; green, 121; blue, 243 }  ,draw opacity=1 ][line width=1.5]  (223.23,97) .. controls (223.23,61.28) and (253.46,32.33) .. (290.75,32.33) .. controls (328.04,32.33) and (358.27,61.28) .. (358.27,97) .. controls (358.27,132.72) and (328.04,161.67) .. (290.75,161.67) .. controls (253.46,161.67) and (223.23,132.72) .. (223.23,97) -- cycle ;
%Shape: Ellipse [id:dp06376891741957036] 
\draw  [color={rgb, 255:red, 16; green, 121; blue, 243 }  ,draw opacity=1 ][line width=1.5]  (249.76,98.85) .. controls (249.76,77.16) and (268.11,59.58) .. (290.75,59.58) .. controls (313.39,59.58) and (331.74,77.16) .. (331.74,98.85) .. controls (331.74,120.53) and (313.39,138.11) .. (290.75,138.11) .. controls (268.11,138.11) and (249.76,120.53) .. (249.76,98.85) -- cycle ;

% Text Node
\draw (288.82,13.05) node [anchor=north west][inner sep=0.75pt]  [color={rgb, 255:red, 208; green, 2; blue, 27 }  ,opacity=1 ]  {$A$};
% Text Node
\draw (287.9,173.81) node [anchor=north west][inner sep=0.75pt]  [color={rgb, 255:red, 208; green, 2; blue, 27 }  ,opacity=1 ]  {$B$};
% Text Node
\draw (230.96,92.51) node [anchor=north west][inner sep=0.75pt]  [color={rgb, 255:red, 208; green, 2; blue, 27 }  ,opacity=1 ]  {$C$};
% Text Node
\draw (273.38,104.52) node [anchor=north west][inner sep=0.75pt]  [color={rgb, 255:red, 208; green, 2; blue, 27 }  ,opacity=1 ]  {$D$};
% Text Node
\draw (340.96,93.43) node [anchor=north west][inner sep=0.75pt]  [color={rgb, 255:red, 208; green, 2; blue, 27 }  ,opacity=1 ]  {$E$};
% Text Node
\draw (123.17,5.87) node [anchor=north west][inner sep=0.75pt]    {$a_{1}$};
% Text Node
\draw (104.56,45.44) node [anchor=north west][inner sep=0.75pt]    {$a_{2}$};
% Text Node
\draw (78.51,61.14) node [anchor=north west][inner sep=0.75pt]    {$a_{3}$};
% Text Node
\draw (122.88,64.84) node [anchor=north west][inner sep=0.75pt]    {$a_{4}$};
% Text Node
\draw (47.65,115.65) node [anchor=north west][inner sep=0.75pt]    {$a_{5}$};
% Text Node
\draw (97.81,126.74) node [anchor=north west][inner sep=0.75pt]    {$a_{6}$};
% Text Node
\draw (69.83,135.05) node [anchor=north west][inner sep=0.75pt]    {$a_{7}$};
% Text Node
\draw (52.47,179.4) node [anchor=north west][inner sep=0.75pt]    {$a_{8}$};
% Text Node
\draw (52.53,208.56) node [anchor=north west][inner sep=0.75pt]  [font=\small] [align=left] {Grafo G};
% Text Node
\draw (77.63,24.14) node [anchor=north west][inner sep=0.75pt]  [color={rgb, 255:red, 208; green, 2; blue, 27 }  ,opacity=1 ]  {$1$};
% Text Node
\draw (161.54,33.38) node [anchor=north west][inner sep=0.75pt]  [color={rgb, 255:red, 208; green, 2; blue, 27 }  ,opacity=1 ]  {$2$};
% Text Node
\draw (77.63,87.89) node [anchor=north west][inner sep=0.75pt]  [color={rgb, 255:red, 208; green, 2; blue, 27 }  ,opacity=1 ]  {$3$};
% Text Node
\draw (19.75,153.48) node [anchor=north west][inner sep=0.75pt]  [color={rgb, 255:red, 208; green, 2; blue, 27 }  ,opacity=1 ]  {$4$};
% Text Node
\draw (101.74,166.42) node [anchor=north west][inner sep=0.75pt]  [color={rgb, 255:red, 208; green, 2; blue, 27 }  ,opacity=1 ]  {$5$};
% Text Node
\draw (252.43,56.68) node [anchor=north west][inner sep=0.75pt]    {$e_{1}$};
% Text Node
\draw (268.54,81.47) node [anchor=north west][inner sep=0.75pt]    {$e_{2}$};
% Text Node
\draw (305.2,83.32) node [anchor=north west][inner sep=0.75pt]    {$e_{3}$};
% Text Node
\draw (319.66,124.89) node [anchor=north west][inner sep=0.75pt]    {$e_{4}$};
% Text Node
\draw (211.63,112.88) node [anchor=north west][inner sep=0.75pt]    {$e_{5}$};
% Text Node
\draw (353.42,47.29) node [anchor=north west][inner sep=0.75pt]    {$e_{6}$};
% Text Node
\draw (294.4,42.18) node [anchor=north west][inner sep=0.75pt]    {$e_{7}$};
% Text Node
\draw (293.43,137.06) node [anchor=north west][inner sep=0.75pt]    {$e_{8}$};
% Text Node
\draw (268.53,207.56) node [anchor=north west][inner sep=0.75pt]  [font=\small] [align=left] {Grafo H};


\end{tikzpicture}

\end{ftst}


%=====

\begin{ftst}{Exercícios}{Tipos de Grafos}

Prove que os dois grafos a seguir não são isomorfos.

\vone
\centering


\tikzset{every picture/.style={line width=0.75pt}} %set default line width to 0.75pt        

\begin{tikzpicture}[x=0.75pt,y=0.75pt,yscale=-1,xscale=1]
%uncomment if require: \path (0,155); %set diagram left start at 0, and has height of 155

%Shape: Ellipse [id:dp12536594919516397] 
\draw  [draw opacity=0][fill={rgb, 255:red, 15; green, 118; blue, 237 }  ,fill opacity=1 ] (18.55,26.91) .. controls (21.33,24.36) and (25.86,24.42) .. (28.67,27.06) .. controls (31.47,29.7) and (31.48,33.9) .. (28.7,36.46) .. controls (25.91,39.01) and (21.38,38.94) .. (18.58,36.31) .. controls (15.78,33.67) and (15.76,29.46) .. (18.55,26.91) -- cycle ;
%Shape: Ellipse [id:dp4126233399315966] 
\draw  [draw opacity=0][fill={rgb, 255:red, 15; green, 118; blue, 237 }  ,fill opacity=1 ] (18.55,92.11) .. controls (21.33,89.55) and (25.86,89.62) .. (28.67,92.26) .. controls (31.47,94.89) and (31.48,99.1) .. (28.7,101.66) .. controls (25.91,104.21) and (21.38,104.14) .. (18.58,101.51) .. controls (15.78,98.87) and (15.76,94.66) .. (18.55,92.11) -- cycle ;
%Straight Lines [id:da6803406363410178] 
\draw [color={rgb, 255:red, 16; green, 121; blue, 243 }  ,draw opacity=1 ][fill={rgb, 255:red, 0; green, 0; blue, 0 }  ,fill opacity=1 ][line width=1.5]    (23.62,31.68) -- (23.62,96.88) ;
%Shape: Ellipse [id:dp39901848215427194] 
\draw  [draw opacity=0][fill={rgb, 255:red, 15; green, 118; blue, 237 }  ,fill opacity=1 ] (83.87,27.82) .. controls (86.66,25.26) and (91.19,25.33) .. (93.99,27.96) .. controls (96.8,30.6) and (96.81,34.81) .. (94.03,37.36) .. controls (91.24,39.92) and (86.71,39.85) .. (83.91,37.21) .. controls (81.1,34.58) and (81.09,30.37) .. (83.87,27.82) -- cycle ;
%Shape: Ellipse [id:dp22255253091413918] 
\draw  [draw opacity=0][fill={rgb, 255:red, 15; green, 118; blue, 237 }  ,fill opacity=1 ] (83.87,93.01) .. controls (86.66,90.46) and (91.19,90.53) .. (93.99,93.16) .. controls (96.8,95.8) and (96.81,100.01) .. (94.03,102.56) .. controls (91.24,105.12) and (86.71,105.05) .. (83.91,102.41) .. controls (81.1,99.78) and (81.09,95.57) .. (83.87,93.01) -- cycle ;
%Straight Lines [id:da6055452552865026] 
\draw [color={rgb, 255:red, 15; green, 118; blue, 237 }  ,draw opacity=1 ][fill={rgb, 255:red, 0; green, 0; blue, 0 }  ,fill opacity=1 ][line width=1.5]    (88.95,32.59) -- (88.95,97.79) ;
%Straight Lines [id:da7864970127723308] 
\draw [color={rgb, 255:red, 15; green, 118; blue, 237 }  ,draw opacity=1 ][fill={rgb, 255:red, 0; green, 0; blue, 0 }  ,fill opacity=1 ][line width=1.5]    (23.62,31.68) -- (88.95,32.59) ;
%Straight Lines [id:da7133886850656368] 
\draw [color={rgb, 255:red, 15; green, 118; blue, 237 }  ,draw opacity=1 ][fill={rgb, 255:red, 0; green, 0; blue, 0 }  ,fill opacity=1 ][line width=1.5]    (23.62,96.88) -- (88.95,97.79) ;
%Straight Lines [id:da5185915981888158] 
\draw [color={rgb, 255:red, 15; green, 118; blue, 237 }  ,draw opacity=1 ][fill={rgb, 255:red, 0; green, 0; blue, 0 }  ,fill opacity=1 ][line width=1.5]    (88.95,32.59) -- (23.62,96.88) ;
%Straight Lines [id:da19696273077754722] 
\draw [color={rgb, 255:red, 15; green, 118; blue, 237 }  ,draw opacity=1 ][fill={rgb, 255:red, 0; green, 0; blue, 0 }  ,fill opacity=1 ][line width=1.5]    (23.62,31.68) -- (88.95,97.79) ;
%Shape: Ellipse [id:dp5067299625850306] 
\draw  [draw opacity=0][fill={rgb, 255:red, 15; green, 118; blue, 237 }  ,fill opacity=1 ] (253.55,27.91) .. controls (256.33,25.36) and (260.86,25.42) .. (263.67,28.06) .. controls (266.47,30.7) and (266.48,34.9) .. (263.7,37.46) .. controls (260.91,40.01) and (256.38,39.94) .. (253.58,37.31) .. controls (250.78,34.67) and (250.76,30.46) .. (253.55,27.91) -- cycle ;
%Shape: Ellipse [id:dp5404576643204713] 
\draw  [draw opacity=0][fill={rgb, 255:red, 15; green, 118; blue, 237 }  ,fill opacity=1 ] (253.55,93.11) .. controls (256.33,90.55) and (260.86,90.62) .. (263.67,93.26) .. controls (266.47,95.89) and (266.48,100.1) .. (263.7,102.66) .. controls (260.91,105.21) and (256.38,105.14) .. (253.58,102.51) .. controls (250.78,99.87) and (250.76,95.66) .. (253.55,93.11) -- cycle ;
%Straight Lines [id:da2912843532175886] 
\draw [color={rgb, 255:red, 16; green, 121; blue, 243 }  ,draw opacity=1 ][fill={rgb, 255:red, 0; green, 0; blue, 0 }  ,fill opacity=1 ][line width=1.5]    (258.62,32.68) -- (258.62,97.88) ;
%Shape: Ellipse [id:dp3406577818091927] 
\draw  [draw opacity=0][fill={rgb, 255:red, 15; green, 118; blue, 237 }  ,fill opacity=1 ] (318.87,28.82) .. controls (321.66,26.26) and (326.19,26.33) .. (328.99,28.96) .. controls (331.8,31.6) and (331.81,35.81) .. (329.03,38.36) .. controls (326.24,40.92) and (321.71,40.85) .. (318.91,38.21) .. controls (316.1,35.58) and (316.09,31.37) .. (318.87,28.82) -- cycle ;
%Shape: Ellipse [id:dp17704059973450392] 
\draw  [draw opacity=0][fill={rgb, 255:red, 15; green, 118; blue, 237 }  ,fill opacity=1 ] (318.87,94.01) .. controls (321.66,91.46) and (326.19,91.53) .. (328.99,94.16) .. controls (331.8,96.8) and (331.81,101.01) .. (329.03,103.56) .. controls (326.24,106.12) and (321.71,106.05) .. (318.91,103.41) .. controls (316.1,100.78) and (316.09,96.57) .. (318.87,94.01) -- cycle ;
%Straight Lines [id:da020826875790622035] 
\draw [color={rgb, 255:red, 15; green, 118; blue, 237 }  ,draw opacity=1 ][fill={rgb, 255:red, 0; green, 0; blue, 0 }  ,fill opacity=1 ][line width=1.5]    (323.95,33.59) -- (323.95,98.79) ;
%Straight Lines [id:da5797162933417743] 
\draw [color={rgb, 255:red, 15; green, 118; blue, 237 }  ,draw opacity=1 ][fill={rgb, 255:red, 0; green, 0; blue, 0 }  ,fill opacity=1 ][line width=1.5]    (258.62,32.68) -- (323.95,33.59) ;
%Straight Lines [id:da5789121412953988] 
\draw [color={rgb, 255:red, 15; green, 118; blue, 237 }  ,draw opacity=1 ][fill={rgb, 255:red, 0; green, 0; blue, 0 }  ,fill opacity=1 ][line width=1.5]    (258.62,97.88) -- (323.95,98.79) ;
%Straight Lines [id:da9400740791814897] 
\draw [color={rgb, 255:red, 15; green, 118; blue, 237 }  ,draw opacity=1 ][fill={rgb, 255:red, 0; green, 0; blue, 0 }  ,fill opacity=1 ][line width=1.5]    (258.62,32.68) -- (323.95,98.79) ;
%Curve Lines [id:da043007482947356346] 
\draw [color={rgb, 255:red, 16; green, 121; blue, 243 }  ,draw opacity=1 ][line width=1.5]    (258.62,32.68) .. controls (308.11,-23.96) and (387.11,34.04) .. (323.95,98.79) ;

% Text Node
\draw (33.53,119.56) node [anchor=north west][inner sep=0.75pt]  [font=\small] [align=left] {Grafo G};
% Text Node
\draw (268.53,120.56) node [anchor=north west][inner sep=0.75pt]  [font=\small] [align=left] {Grafo G};


\end{tikzpicture}

\end{ftst}


%=====

\begin{ftst}{Grafos Planares}{Tipos de Grafos}

Um grafo planar é um grafo que pode ser representado em um plano de modo que suas arestas não se interceptam.

\vone
\centering


\tikzset{every picture/.style={line width=0.75pt}} %set default line width to 0.75pt        

\begin{tikzpicture}[x=0.75pt,y=0.75pt,yscale=-1,xscale=1]
%uncomment if require: \path (0,163); %set diagram left start at 0, and has height of 163

%Shape: Ellipse [id:dp7748248031867107] 
\draw  [draw opacity=0][fill={rgb, 255:red, 15; green, 118; blue, 237 }  ,fill opacity=1 ] (76.4,31.3) .. controls (80.18,31.47) and (83.33,34.72) .. (83.45,38.57) .. controls (83.56,42.42) and (80.59,45.4) .. (76.82,45.23) .. controls (73.04,45.07) and (69.89,41.81) .. (69.77,37.96) .. controls (69.66,34.12) and (72.63,31.13) .. (76.4,31.3) -- cycle ;
%Shape: Ellipse [id:dp5270384232513423] 
\draw  [draw opacity=0][fill={rgb, 255:red, 15; green, 118; blue, 237 }  ,fill opacity=1 ] (30.25,77.36) .. controls (34.03,77.52) and (37.18,80.78) .. (37.29,84.63) .. controls (37.41,88.47) and (34.44,91.45) .. (30.67,91.29) .. controls (26.89,91.12) and (23.74,87.87) .. (23.62,84.02) .. controls (23.51,80.17) and (26.48,77.19) .. (30.25,77.36) -- cycle ;
%Straight Lines [id:da41663362556157524] 
\draw [color={rgb, 255:red, 16; green, 121; blue, 243 }  ,draw opacity=1 ][fill={rgb, 255:red, 0; green, 0; blue, 0 }  ,fill opacity=1 ][line width=1.5]    (76.61,38.27) -- (30.46,84.32) ;
%Shape: Ellipse [id:dp307378844800797] 
\draw  [draw opacity=0][fill={rgb, 255:red, 15; green, 118; blue, 237 }  ,fill opacity=1 ] (121.91,78.18) .. controls (125.68,78.35) and (128.84,81.61) .. (128.95,85.45) .. controls (129.06,89.3) and (126.1,92.28) .. (122.32,92.11) .. controls (118.55,91.95) and (115.39,88.69) .. (115.28,84.85) .. controls (115.16,81) and (118.13,78.02) .. (121.91,78.18) -- cycle ;
%Shape: Ellipse [id:dp49363140495247904] 
\draw  [draw opacity=0][fill={rgb, 255:red, 15; green, 118; blue, 237 }  ,fill opacity=1 ] (75.76,124.24) .. controls (79.53,124.41) and (82.68,127.66) .. (82.8,131.51) .. controls (82.91,135.35) and (79.94,138.34) .. (76.17,138.17) .. controls (72.39,138) and (69.24,134.75) .. (69.13,130.9) .. controls (69.01,127.05) and (71.98,124.07) .. (75.76,124.24) -- cycle ;
%Straight Lines [id:da2138034134787543] 
\draw [color={rgb, 255:red, 15; green, 118; blue, 237 }  ,draw opacity=1 ][fill={rgb, 255:red, 0; green, 0; blue, 0 }  ,fill opacity=1 ][line width=1.5]    (122.11,85.15) -- (75.96,131.2) ;
%Straight Lines [id:da9687693942761628] 
\draw [color={rgb, 255:red, 15; green, 118; blue, 237 }  ,draw opacity=1 ][fill={rgb, 255:red, 0; green, 0; blue, 0 }  ,fill opacity=1 ][line width=1.5]    (76.61,38.27) -- (122.11,85.15) ;
%Straight Lines [id:da2697679454967321] 
\draw [color={rgb, 255:red, 15; green, 118; blue, 237 }  ,draw opacity=1 ][fill={rgb, 255:red, 0; green, 0; blue, 0 }  ,fill opacity=1 ][line width=1.5]    (30.46,84.32) -- (75.96,131.2) ;
%Straight Lines [id:da7120580725748549] 
\draw [color={rgb, 255:red, 15; green, 118; blue, 237 }  ,draw opacity=1 ][fill={rgb, 255:red, 0; green, 0; blue, 0 }  ,fill opacity=1 ][line width=1.5]    (122.11,85.15) -- (30.46,84.32) ;
%Straight Lines [id:da6420267033924327] 
\draw [color={rgb, 255:red, 15; green, 118; blue, 237 }  ,draw opacity=1 ][fill={rgb, 255:red, 0; green, 0; blue, 0 }  ,fill opacity=1 ][line width=1.5]    (76.61,38.27) -- (75.96,131.2) ;
%Curve Lines [id:da8183204265459638] 
\draw [color={rgb, 255:red, 15; green, 118; blue, 237 }  ,draw opacity=1 ][line width=1.5]    (30.46,84.32) .. controls (10.11,67.06) and (34.11,20.06) .. (76.61,38.27) ;
%Curve Lines [id:da8327921492017605] 
\draw [color={rgb, 255:red, 15; green, 118; blue, 237 }  ,draw opacity=1 ][line width=1.5]    (75.96,131.2) .. controls (29.11,151.06) and (14.11,95.06) .. (30.46,84.32) ;
%Shape: Ellipse [id:dp8520257818599006] 
\draw  [draw opacity=0][fill={rgb, 255:red, 15; green, 118; blue, 237 }  ,fill opacity=1 ] (291.4,32.3) .. controls (295.18,32.47) and (298.33,35.72) .. (298.45,39.57) .. controls (298.56,43.42) and (295.59,46.4) .. (291.82,46.23) .. controls (288.04,46.07) and (284.89,42.81) .. (284.77,38.96) .. controls (284.66,35.12) and (287.63,32.13) .. (291.4,32.3) -- cycle ;
%Shape: Ellipse [id:dp22912846854362856] 
\draw  [draw opacity=0][fill={rgb, 255:red, 15; green, 118; blue, 237 }  ,fill opacity=1 ] (245.25,78.36) .. controls (249.03,78.52) and (252.18,81.78) .. (252.29,85.63) .. controls (252.41,89.47) and (249.44,92.45) .. (245.67,92.29) .. controls (241.89,92.12) and (238.74,88.87) .. (238.62,85.02) .. controls (238.51,81.17) and (241.48,78.19) .. (245.25,78.36) -- cycle ;
%Straight Lines [id:da08750807811322314] 
\draw [color={rgb, 255:red, 16; green, 121; blue, 243 }  ,draw opacity=1 ][fill={rgb, 255:red, 0; green, 0; blue, 0 }  ,fill opacity=1 ][line width=1.5]    (291.61,39.27) -- (245.46,85.32) ;
%Shape: Ellipse [id:dp11479439664611935] 
\draw  [draw opacity=0][fill={rgb, 255:red, 15; green, 118; blue, 237 }  ,fill opacity=1 ] (336.91,79.18) .. controls (340.68,79.35) and (343.84,82.61) .. (343.95,86.45) .. controls (344.06,90.3) and (341.1,93.28) .. (337.32,93.11) .. controls (333.55,92.95) and (330.39,89.69) .. (330.28,85.85) .. controls (330.16,82) and (333.13,79.02) .. (336.91,79.18) -- cycle ;
%Shape: Ellipse [id:dp6283334101894364] 
\draw  [draw opacity=0][fill={rgb, 255:red, 15; green, 118; blue, 237 }  ,fill opacity=1 ] (290.76,125.24) .. controls (294.53,125.41) and (297.68,128.66) .. (297.8,132.51) .. controls (297.91,136.35) and (294.94,139.34) .. (291.17,139.17) .. controls (287.39,139) and (284.24,135.75) .. (284.13,131.9) .. controls (284.01,128.05) and (286.98,125.07) .. (290.76,125.24) -- cycle ;
%Straight Lines [id:da7409963607950782] 
\draw [color={rgb, 255:red, 15; green, 118; blue, 237 }  ,draw opacity=1 ][fill={rgb, 255:red, 0; green, 0; blue, 0 }  ,fill opacity=1 ][line width=1.5]    (337.11,86.15) -- (290.96,132.2) ;
%Straight Lines [id:da08331299212213028] 
\draw [color={rgb, 255:red, 15; green, 118; blue, 237 }  ,draw opacity=1 ][fill={rgb, 255:red, 0; green, 0; blue, 0 }  ,fill opacity=1 ][line width=1.5]    (291.61,39.27) -- (337.11,86.15) ;
%Straight Lines [id:da4128994649916242] 
\draw [color={rgb, 255:red, 15; green, 118; blue, 237 }  ,draw opacity=1 ][fill={rgb, 255:red, 0; green, 0; blue, 0 }  ,fill opacity=1 ][line width=1.5]    (245.46,85.32) -- (290.96,132.2) ;
%Straight Lines [id:da5509314471531617] 
\draw [color={rgb, 255:red, 15; green, 118; blue, 237 }  ,draw opacity=1 ][fill={rgb, 255:red, 0; green, 0; blue, 0 }  ,fill opacity=1 ][line width=1.5]    (337.11,86.15) -- (245.46,85.32) ;
%Curve Lines [id:da4358243239346997] 
\draw [color={rgb, 255:red, 15; green, 118; blue, 237 }  ,draw opacity=1 ][line width=1.5]    (245.46,85.32) .. controls (225.11,68.06) and (249.11,21.06) .. (291.61,39.27) ;
%Curve Lines [id:da508099961137481] 
\draw [color={rgb, 255:red, 15; green, 118; blue, 237 }  ,draw opacity=1 ][line width=1.5]    (290.96,132.2) .. controls (244.11,152.06) and (229.11,96.06) .. (245.46,85.32) ;
%Curve Lines [id:da8885767908282762] 
\draw [color={rgb, 255:red, 15; green, 118; blue, 237 }  ,draw opacity=1 ][line width=1.5]    (291.61,39.27) .. controls (409.11,45.06) and (393.11,133.06) .. (290.96,132.2) ;
%Right Arrow [id:dp43191010002441454] 
\draw  [draw opacity=0][fill={rgb, 255:red, 245; green, 166; blue, 35 }  ,fill opacity=1 ] (170,76.3) -- (192.56,76.3) -- (192.56,70.06) -- (207.6,82.53) -- (192.56,95) -- (192.56,88.77) -- (170,88.77) -- cycle ;




\end{tikzpicture}

\end{ftst}


%=====

\begin{ftst}{Grafos Planares}{Tipos de Grafos}
\begin{itemize}
    \item A representação planar de um grafo divide o plano em regiões (uma é ilimitada).
    \item Seja $G$ um grafo plano conexo com $|E|$ arestas e $|V|$ vértices, com um número $R$ de regiões na representação planar de $G$, temos: $R=|E| - |V|+2$ (fórmula de Euler)
\end{itemize}

\vone
\centering


\tikzset{every picture/.style={line width=0.75pt}} %set default line width to 0.75pt        

\begin{tikzpicture}[x=0.75pt,y=0.75pt,yscale=-1,xscale=1]
%uncomment if require: \path (0,151); %set diagram left start at 0, and has height of 151

%Shape: Ellipse [id:dp6424091519689461] 
\draw  [draw opacity=0][fill={rgb, 255:red, 15; green, 118; blue, 237 }  ,fill opacity=1 ] (108.4,27.3) .. controls (112.18,27.47) and (115.33,30.72) .. (115.45,34.57) .. controls (115.56,38.42) and (112.59,41.4) .. (108.82,41.23) .. controls (105.04,41.07) and (101.89,37.81) .. (101.77,33.96) .. controls (101.66,30.12) and (104.63,27.13) .. (108.4,27.3) -- cycle ;
%Shape: Ellipse [id:dp4868952713488026] 
\draw  [draw opacity=0][fill={rgb, 255:red, 15; green, 118; blue, 237 }  ,fill opacity=1 ] (62.25,73.36) .. controls (66.03,73.52) and (69.18,76.78) .. (69.29,80.63) .. controls (69.41,84.47) and (66.44,87.45) .. (62.67,87.29) .. controls (58.89,87.12) and (55.74,83.87) .. (55.62,80.02) .. controls (55.51,76.17) and (58.48,73.19) .. (62.25,73.36) -- cycle ;
%Straight Lines [id:da44467552426780155] 
\draw [color={rgb, 255:red, 16; green, 121; blue, 243 }  ,draw opacity=1 ][fill={rgb, 255:red, 0; green, 0; blue, 0 }  ,fill opacity=1 ][line width=1.5]    (108.61,34.27) -- (62.46,80.32) ;
%Shape: Ellipse [id:dp3390874097767833] 
\draw  [draw opacity=0][fill={rgb, 255:red, 15; green, 118; blue, 237 }  ,fill opacity=1 ] (153.91,74.18) .. controls (157.68,74.35) and (160.84,77.61) .. (160.95,81.45) .. controls (161.06,85.3) and (158.1,88.28) .. (154.32,88.11) .. controls (150.55,87.95) and (147.39,84.69) .. (147.28,80.85) .. controls (147.16,77) and (150.13,74.02) .. (153.91,74.18) -- cycle ;
%Shape: Ellipse [id:dp5106914064389432] 
\draw  [draw opacity=0][fill={rgb, 255:red, 15; green, 118; blue, 237 }  ,fill opacity=1 ] (107.76,120.24) .. controls (111.53,120.41) and (114.68,123.66) .. (114.8,127.51) .. controls (114.91,131.35) and (111.94,134.34) .. (108.17,134.17) .. controls (104.39,134) and (101.24,130.75) .. (101.13,126.9) .. controls (101.01,123.05) and (103.98,120.07) .. (107.76,120.24) -- cycle ;
%Straight Lines [id:da763734854884363] 
\draw [color={rgb, 255:red, 15; green, 118; blue, 237 }  ,draw opacity=1 ][fill={rgb, 255:red, 0; green, 0; blue, 0 }  ,fill opacity=1 ][line width=1.5]    (154.11,81.15) -- (107.96,127.2) ;
%Straight Lines [id:da588551383268862] 
\draw [color={rgb, 255:red, 15; green, 118; blue, 237 }  ,draw opacity=1 ][fill={rgb, 255:red, 0; green, 0; blue, 0 }  ,fill opacity=1 ][line width=1.5]    (108.61,34.27) -- (154.11,81.15) ;
%Straight Lines [id:da14849799274214592] 
\draw [color={rgb, 255:red, 15; green, 118; blue, 237 }  ,draw opacity=1 ][fill={rgb, 255:red, 0; green, 0; blue, 0 }  ,fill opacity=1 ][line width=1.5]    (62.46,80.32) -- (107.96,127.2) ;
%Straight Lines [id:da7402882637314507] 
\draw [color={rgb, 255:red, 15; green, 118; blue, 237 }  ,draw opacity=1 ][fill={rgb, 255:red, 0; green, 0; blue, 0 }  ,fill opacity=1 ][line width=1.5]    (154.11,81.15) -- (62.46,80.32) ;
%Curve Lines [id:da2331445702439543] 
\draw [color={rgb, 255:red, 15; green, 118; blue, 237 }  ,draw opacity=1 ][line width=1.5]    (108.61,34.27) .. controls (226.11,40.06) and (210.11,128.06) .. (107.96,127.2) ;

% Text Node
\draw (32,35.4) node [anchor=north west][inner sep=0.75pt]  [color={rgb, 255:red, 245; green, 166; blue, 35 }  ,opacity=1 ]  {$R_{4}$};
% Text Node
\draw (102,90.4) node [anchor=north west][inner sep=0.75pt]  [color={rgb, 255:red, 245; green, 166; blue, 35 }  ,opacity=1 ]  {$R_{2}$};
% Text Node
\draw (166,74.4) node [anchor=north west][inner sep=0.75pt]  [color={rgb, 255:red, 245; green, 166; blue, 35 }  ,opacity=1 ]  {$R_{3}$};
% Text Node
\draw (101,55.4) node [anchor=north west][inner sep=0.75pt]  [color={rgb, 255:red, 245; green, 166; blue, 35 }  ,opacity=1 ]  {$R_{1}$};
% Text Node
\draw (223,59.4) node [anchor=north west][inner sep=0.75pt]  [color={rgb, 255:red, 0; green, 0; blue, 0 }  ,opacity=1 ]  {$ \begin{array}{l}
R\ =\ 6\ -\ 4\ +\ 2\\
R\ =\ 4
\end{array}$};


\end{tikzpicture}

\end{ftst}


%=====

\begin{ftst}{Grafos Planares}{Tipos de Grafos}

Seja $G$ é um grafo planar simples com $|V|$ vértices e $|E|$ arestas, tal que $|V| \geq 3$. Temos que: $|E| \leq 3|V| - 6$.

\vone
\vone
\vone
\vone

\centering


\tikzset{every picture/.style={line width=0.75pt}} %set default line width to 0.75pt        

\begin{tikzpicture}[x=0.75pt,y=0.75pt,yscale=-1,xscale=1]
%uncomment if require: \path (0,129); %set diagram left start at 0, and has height of 129

%Shape: Ellipse [id:dp8824645259000807] 
\draw  [draw opacity=0][fill={rgb, 255:red, 15; green, 118; blue, 237 }  ,fill opacity=1 ] (83.4,5.3) .. controls (87.18,5.47) and (90.33,8.72) .. (90.45,12.57) .. controls (90.56,16.42) and (87.59,19.4) .. (83.82,19.23) .. controls (80.04,19.07) and (76.89,15.81) .. (76.77,11.96) .. controls (76.66,8.12) and (79.63,5.13) .. (83.4,5.3) -- cycle ;
%Shape: Ellipse [id:dp19463839491212043] 
\draw  [draw opacity=0][fill={rgb, 255:red, 15; green, 118; blue, 237 }  ,fill opacity=1 ] (37.25,51.36) .. controls (41.03,51.52) and (44.18,54.78) .. (44.29,58.63) .. controls (44.41,62.47) and (41.44,65.45) .. (37.67,65.29) .. controls (33.89,65.12) and (30.74,61.87) .. (30.62,58.02) .. controls (30.51,54.17) and (33.48,51.19) .. (37.25,51.36) -- cycle ;
%Straight Lines [id:da14445363951451218] 
\draw [color={rgb, 255:red, 16; green, 121; blue, 243 }  ,draw opacity=1 ][fill={rgb, 255:red, 0; green, 0; blue, 0 }  ,fill opacity=1 ][line width=1.5]    (83.61,12.27) -- (37.46,58.32) ;
%Shape: Ellipse [id:dp008236776664483747] 
\draw  [draw opacity=0][fill={rgb, 255:red, 15; green, 118; blue, 237 }  ,fill opacity=1 ] (128.91,52.18) .. controls (132.68,52.35) and (135.84,55.61) .. (135.95,59.45) .. controls (136.06,63.3) and (133.1,66.28) .. (129.32,66.11) .. controls (125.55,65.95) and (122.39,62.69) .. (122.28,58.85) .. controls (122.16,55) and (125.13,52.02) .. (128.91,52.18) -- cycle ;
%Shape: Ellipse [id:dp7929686215334855] 
\draw  [draw opacity=0][fill={rgb, 255:red, 15; green, 118; blue, 237 }  ,fill opacity=1 ] (82.76,98.24) .. controls (86.53,98.41) and (89.68,101.66) .. (89.8,105.51) .. controls (89.91,109.35) and (86.94,112.34) .. (83.17,112.17) .. controls (79.39,112) and (76.24,108.75) .. (76.13,104.9) .. controls (76.01,101.05) and (78.98,98.07) .. (82.76,98.24) -- cycle ;
%Straight Lines [id:da24639820678501922] 
\draw [color={rgb, 255:red, 15; green, 118; blue, 237 }  ,draw opacity=1 ][fill={rgb, 255:red, 0; green, 0; blue, 0 }  ,fill opacity=1 ][line width=1.5]    (129.11,59.15) -- (82.96,105.2) ;
%Straight Lines [id:da9492046865987103] 
\draw [color={rgb, 255:red, 15; green, 118; blue, 237 }  ,draw opacity=1 ][fill={rgb, 255:red, 0; green, 0; blue, 0 }  ,fill opacity=1 ][line width=1.5]    (83.61,12.27) -- (129.11,59.15) ;
%Straight Lines [id:da1986454721607971] 
\draw [color={rgb, 255:red, 15; green, 118; blue, 237 }  ,draw opacity=1 ][fill={rgb, 255:red, 0; green, 0; blue, 0 }  ,fill opacity=1 ][line width=1.5]    (37.46,58.32) -- (82.96,105.2) ;
%Straight Lines [id:da17378448524728318] 
\draw [color={rgb, 255:red, 15; green, 118; blue, 237 }  ,draw opacity=1 ][fill={rgb, 255:red, 0; green, 0; blue, 0 }  ,fill opacity=1 ][line width=1.5]    (129.11,59.15) -- (37.46,58.32) ;
%Curve Lines [id:da13127401435563257] 
\draw [color={rgb, 255:red, 15; green, 118; blue, 237 }  ,draw opacity=1 ][line width=1.5]    (83.61,12.27) .. controls (201.11,18.06) and (185.11,106.06) .. (82.96,105.2) ;

% Text Node
\draw (7,13.4) node [anchor=north west][inner sep=0.75pt]  [color={rgb, 255:red, 245; green, 166; blue, 35 }  ,opacity=1 ]  {$R_{4}$};
% Text Node
\draw (77,68.4) node [anchor=north west][inner sep=0.75pt]  [color={rgb, 255:red, 245; green, 166; blue, 35 }  ,opacity=1 ]  {$R_{2}$};
% Text Node
\draw (141,52.4) node [anchor=north west][inner sep=0.75pt]  [color={rgb, 255:red, 245; green, 166; blue, 35 }  ,opacity=1 ]  {$R_{3}$};
% Text Node
\draw (76,33.4) node [anchor=north west][inner sep=0.75pt]  [color={rgb, 255:red, 245; green, 166; blue, 35 }  ,opacity=1 ]  {$R_{1}$};
% Text Node
\draw (198,37.4) node [anchor=north west][inner sep=0.75pt]  [color={rgb, 255:red, 0; green, 0; blue, 0 }  ,opacity=1 ]  {$ \begin{array}{l}
|E|\ \leqslant \ 3\ *\ 4\ -\ 6\\
|E|\ \leqslant \ 6\\
6\ \leqslant \ 6\\
\end{array}$};


\end{tikzpicture}

\end{ftst}


%=====

\begin{ftst}{Grafos Planares}{Tipos de Grafos}
\begin{itemize}
    \item Todo grafo planar simples tem um vértice de grau no máximo 5.
    \item Os grafos completos $K_5$ e $K{3,3}$ não são planares.
\end{itemize}
\vone
\vone
\centering


\tikzset{every picture/.style={line width=0.75pt}} %set default line width to 0.75pt        

\begin{tikzpicture}[x=0.75pt,y=0.75pt,yscale=-1,xscale=1]
%uncomment if require: \path (0,161); %set diagram left start at 0, and has height of 161

%Shape: Circle [id:dp742266466022518] 
\draw  [draw opacity=0][fill={rgb, 255:red, 15; green, 118; blue, 237 }  ,fill opacity=1 ] (64,31.81) .. controls (64,27.49) and (67.49,24) .. (71.81,24) .. controls (76.12,24) and (79.61,27.49) .. (79.61,31.81) .. controls (79.61,36.12) and (76.12,39.61) .. (71.81,39.61) .. controls (67.49,39.61) and (64,36.12) .. (64,31.81) -- cycle ;
%Shape: Circle [id:dp5902964350745019] 
\draw  [draw opacity=0][fill={rgb, 255:red, 15; green, 118; blue, 237 }  ,fill opacity=1 ] (81,103.81) .. controls (81,99.49) and (84.49,96) .. (88.81,96) .. controls (93.12,96) and (96.61,99.49) .. (96.61,103.81) .. controls (96.61,108.12) and (93.12,111.61) .. (88.81,111.61) .. controls (84.49,111.61) and (81,108.12) .. (81,103.81) -- cycle ;
%Shape: Circle [id:dp9530392712385023] 
\draw  [draw opacity=0][fill={rgb, 255:red, 15; green, 118; blue, 237 }  ,fill opacity=1 ] (99,66.81) .. controls (99,62.49) and (102.49,59) .. (106.81,59) .. controls (111.12,59) and (114.61,62.49) .. (114.61,66.81) .. controls (114.61,71.12) and (111.12,74.61) .. (106.81,74.61) .. controls (102.49,74.61) and (99,71.12) .. (99,66.81) -- cycle ;
%Shape: Circle [id:dp8903023633413498] 
\draw  [draw opacity=0][fill={rgb, 255:red, 15; green, 118; blue, 237 }  ,fill opacity=1 ] (47,102.81) .. controls (47,98.49) and (50.49,95) .. (54.81,95) .. controls (59.12,95) and (62.61,98.49) .. (62.61,102.81) .. controls (62.61,107.12) and (59.12,110.61) .. (54.81,110.61) .. controls (50.49,110.61) and (47,107.12) .. (47,102.81) -- cycle ;
%Shape: Circle [id:dp08086931921463147] 
\draw  [draw opacity=0][fill={rgb, 255:red, 15; green, 118; blue, 237 }  ,fill opacity=1 ] (27,67.81) .. controls (27,63.49) and (30.49,60) .. (34.81,60) .. controls (39.12,60) and (42.61,63.49) .. (42.61,67.81) .. controls (42.61,72.12) and (39.12,75.61) .. (34.81,75.61) .. controls (30.49,75.61) and (27,72.12) .. (27,67.81) -- cycle ;
%Straight Lines [id:da5301075507529325] 
\draw [color={rgb, 255:red, 15; green, 118; blue, 237 }  ,draw opacity=1 ][line width=1.5]    (34.81,67.81) -- (71.81,31.81) ;
%Straight Lines [id:da10853523426024725] 
\draw [color={rgb, 255:red, 15; green, 118; blue, 237 }  ,draw opacity=1 ][line width=1.5]    (106.81,66.81) -- (71.81,31.81) ;
%Straight Lines [id:da04816312350512608] 
\draw [color={rgb, 255:red, 15; green, 118; blue, 237 }  ,draw opacity=1 ][line width=1.5]    (54.81,102.81) -- (34.81,67.81) ;
%Straight Lines [id:da563229743251703] 
\draw [color={rgb, 255:red, 15; green, 118; blue, 237 }  ,draw opacity=1 ][line width=1.5]    (88.81,103.81) -- (106.81,66.81) ;
%Straight Lines [id:da3855013841744166] 
\draw [color={rgb, 255:red, 15; green, 118; blue, 237 }  ,draw opacity=1 ][line width=1.5]    (54.81,102.81) -- (88.81,103.81) ;
%Straight Lines [id:da4439577475289598] 
\draw [color={rgb, 255:red, 15; green, 118; blue, 237 }  ,draw opacity=1 ][line width=1.5]    (54.81,102.81) -- (71.81,31.81) ;
%Straight Lines [id:da8763446748763963] 
\draw [color={rgb, 255:red, 15; green, 118; blue, 237 }  ,draw opacity=1 ][line width=1.5]    (88.81,103.81) -- (34.81,67.81) ;
%Straight Lines [id:da8049967822206365] 
\draw [color={rgb, 255:red, 15; green, 118; blue, 237 }  ,draw opacity=1 ][line width=1.5]    (106.81,66.81) -- (34.81,67.81) ;
%Straight Lines [id:da9979263846231452] 
\draw [color={rgb, 255:red, 15; green, 118; blue, 237 }  ,draw opacity=1 ][line width=1.5]    (88.81,103.81) -- (71.81,31.81) ;
%Straight Lines [id:da9985585398360899] 
\draw [color={rgb, 255:red, 15; green, 118; blue, 237 }  ,draw opacity=1 ][line width=1.5]    (106.81,66.81) -- (54.81,102.81) ;
%Straight Lines [id:da39835957281227685] 
\draw [color={rgb, 255:red, 121; green, 120; blue, 120 }  ,draw opacity=1 ][line width=1.5]    (185.33,20.67) -- (184.33,106.67) ;
%Straight Lines [id:da46629447329855944] 
\draw [color={rgb, 255:red, 121; green, 120; blue, 120 }  ,draw opacity=1 ][line width=1.5]    (185.33,20.67) -- (252.33,106.67) ;
%Straight Lines [id:da886076681519776] 
\draw [color={rgb, 255:red, 121; green, 120; blue, 120 }  ,draw opacity=1 ][line width=1.5]    (185.33,20.67) -- (321.33,106.67) ;
%Straight Lines [id:da8351308406591933] 
\draw [color={rgb, 255:red, 121; green, 120; blue, 120 }  ,draw opacity=1 ][line width=1.5]    (252.33,21.67) -- (184.33,106.67) ;
%Straight Lines [id:da9091762153568039] 
\draw [color={rgb, 255:red, 121; green, 120; blue, 120 }  ,draw opacity=1 ][line width=1.5]    (252.33,21.67) -- (252.33,106.67) ;
%Straight Lines [id:da7337747247298743] 
\draw [color={rgb, 255:red, 121; green, 120; blue, 120 }  ,draw opacity=1 ][line width=1.5]    (252.33,21.67) -- (321.33,106.67) ;
%Straight Lines [id:da33660292872236486] 
\draw [color={rgb, 255:red, 121; green, 120; blue, 120 }  ,draw opacity=1 ][line width=1.5]    (320.33,21.67) -- (184.33,106.67) ;
%Straight Lines [id:da4311715977128252] 
\draw [color={rgb, 255:red, 121; green, 120; blue, 120 }  ,draw opacity=1 ][line width=1.5]    (320.33,21.67) -- (252.33,106.67) ;
%Straight Lines [id:da11561955782961397] 
\draw [color={rgb, 255:red, 121; green, 120; blue, 120 }  ,draw opacity=1 ][line width=1.5]    (320.33,21.67) -- (321.33,106.67) ;
%Shape: Ellipse [id:dp13672353513853253] 
\draw  [draw opacity=0][fill={rgb, 255:red, 245; green, 166; blue, 35 }  ,fill opacity=1 ] (177.92,20.67) .. controls (177.92,16.66) and (181.24,13.41) .. (185.33,13.41) .. controls (189.42,13.41) and (192.74,16.66) .. (192.74,20.67) .. controls (192.74,24.69) and (189.42,27.94) .. (185.33,27.94) .. controls (181.24,27.94) and (177.92,24.69) .. (177.92,20.67) -- cycle ;
%Shape: Ellipse [id:dp1610614475522445] 
\draw  [draw opacity=0][fill={rgb, 255:red, 245; green, 166; blue, 35 }  ,fill opacity=1 ] (244.92,21.67) .. controls (244.92,17.66) and (248.24,14.41) .. (252.33,14.41) .. controls (256.42,14.41) and (259.74,17.66) .. (259.74,21.67) .. controls (259.74,25.69) and (256.42,28.94) .. (252.33,28.94) .. controls (248.24,28.94) and (244.92,25.69) .. (244.92,21.67) -- cycle ;
%Shape: Ellipse [id:dp46414166545397295] 
\draw  [draw opacity=0][fill={rgb, 255:red, 245; green, 166; blue, 35 }  ,fill opacity=1 ] (312.92,21.67) .. controls (312.92,17.66) and (316.24,14.41) .. (320.33,14.41) .. controls (324.42,14.41) and (327.74,17.66) .. (327.74,21.67) .. controls (327.74,25.69) and (324.42,28.94) .. (320.33,28.94) .. controls (316.24,28.94) and (312.92,25.69) .. (312.92,21.67) -- cycle ;
%Shape: Ellipse [id:dp9542768800783743] 
\draw  [draw opacity=0][fill={rgb, 255:red, 15; green, 118; blue, 237 }  ,fill opacity=1 ] (176.92,106.67) .. controls (176.92,102.66) and (180.24,99.41) .. (184.33,99.41) .. controls (188.42,99.41) and (191.74,102.66) .. (191.74,106.67) .. controls (191.74,110.69) and (188.42,113.94) .. (184.33,113.94) .. controls (180.24,113.94) and (176.92,110.69) .. (176.92,106.67) -- cycle ;
%Shape: Ellipse [id:dp07897753579368283] 
\draw  [draw opacity=0][fill={rgb, 255:red, 15; green, 118; blue, 237 }  ,fill opacity=1 ] (244.92,106.67) .. controls (244.92,102.66) and (248.24,99.41) .. (252.33,99.41) .. controls (256.42,99.41) and (259.74,102.66) .. (259.74,106.67) .. controls (259.74,110.69) and (256.42,113.94) .. (252.33,113.94) .. controls (248.24,113.94) and (244.92,110.69) .. (244.92,106.67) -- cycle ;
%Shape: Ellipse [id:dp6237007974814952] 
\draw  [draw opacity=0][fill={rgb, 255:red, 15; green, 118; blue, 237 }  ,fill opacity=1 ] (313.92,106.67) .. controls (313.92,102.66) and (317.24,99.41) .. (321.33,99.41) .. controls (325.42,99.41) and (328.74,102.66) .. (328.74,106.67) .. controls (328.74,110.69) and (325.42,113.94) .. (321.33,113.94) .. controls (317.24,113.94) and (313.92,110.69) .. (313.92,106.67) -- cycle ;

% Text Node
\draw (62,132.4) node [anchor=north west][inner sep=0.75pt]  [color={rgb, 255:red, 0; green, 0; blue, 0 }  ,opacity=1 ]  {$K_{5}$};
% Text Node
\draw (247,129.4) node [anchor=north west][inner sep=0.75pt]  [color={rgb, 255:red, 0; green, 0; blue, 0 }  ,opacity=1 ]  {$K_{\{3,3\}}$};


\end{tikzpicture}

\end{ftst}


%=====


\end{document}